\documentclass[11pt]{article}

\usepackage[letterpaper,margin=0.82in,headheight=15pt]{geometry}
\usepackage[T1]{fontenc}
\usepackage{lmodern}
\usepackage{microtype}
\usepackage{amsmath,amssymb,amsthm,bm,mathtools}
\usepackage{booktabs,longtable,tabularx,array,multirow}
\usepackage{enumitem}
\usepackage{xcolor}
\usepackage{fancyhdr}
\usepackage{graphicx}
\usepackage{tikz-cd}
\usepackage{hyperref}
\usepackage[nameinlink,noabbrev]{cleveref}

\definecolor{navy}{HTML}{17324D}
\definecolor{teal}{HTML}{147D92}
\definecolor{orange}{HTML}{D47A27}
\definecolor{softblue}{HTML}{EEF4F7}
\definecolor{softgray}{HTML}{F3F5F6}
\definecolor{darkgray}{HTML}{46545E}
\definecolor{softorange}{HTML}{FFF4E8}
\definecolor{softgreen}{HTML}{EDF7F0}

\hypersetup{
  colorlinks=true,
  linkcolor=navy,
  citecolor=teal,
  urlcolor=teal,
  pdftitle={Volume VI: Shared Microscopic Inference, Correlated Uncertainty, and Withheld Validation},
  pdfauthor={DeuteronWigner project}
}

\pagestyle{fancy}
\fancyhf{}
\fancyhead[L]{\small\color{darkgray}DeuteronWigner formalism - Volume VI}
\fancyhead[R]{\small\color{darkgray}Inference, uncertainty, and falsifiability}
\fancyfoot[L]{\small\color{darkgray}30 July 2026}
\fancyfoot[R]{\small\color{darkgray}\thepage}
\renewcommand{\headrulewidth}{0.4pt}

\setlist[itemize]{leftmargin=1.5em,itemsep=2pt,topsep=3pt}
\setlist[enumerate]{leftmargin=1.7em,itemsep=2pt,topsep=3pt}
\setlength{\parindent}{0pt}
\setlength{\parskip}{5pt}
\setlength{\emergencystretch}{2em}
\renewcommand{\arraystretch}{1.16}
\allowdisplaybreaks

\newtheorem{definition}{Definition}[section]
\newtheorem{axiom}[definition]{Axiom}
\newtheorem{principle}[definition]{Design principle}
\newtheorem{requirement}[definition]{Requirement}
\newtheorem{proposition}[definition]{Proposition}
\newtheorem{theorem}[definition]{Theorem}
\newtheorem{lemma}[definition]{Lemma}
\newtheorem{corollary}[definition]{Corollary}
\newtheorem{remark}[definition]{Remark}

\crefname{definition}{definition}{definitions}
\crefname{axiom}{axiom}{axioms}
\crefname{principle}{principle}{principles}
\crefname{requirement}{requirement}{requirements}
\crefname{proposition}{proposition}{propositions}
\crefname{theorem}{theorem}{theorems}
\crefname{lemma}{lemma}{lemmas}
\crefname{corollary}{corollary}{corollaries}
\crefname{remark}{remark}{remarks}

\newcommand{\kT}{\bm{k}_{T}}
\newcommand{\qT}{\bm{q}_{T}}
\newcommand{\DT}{\bm{\Delta}_{T}}
\newcommand{\bD}{\bm{b}_{\Delta}}
\newcommand{\bM}{\bm{b}_{\mathrm{TMD}}}
\newcommand{\dd}{\mathrm{d}}
\newcommand{\Tr}{\operatorname{Tr}}
\newcommand{\tr}{\operatorname{tr}}
\newcommand{\diag}{\operatorname{diag}}
\newcommand{\rank}{\operatorname{rank}}
\newcommand{\Span}{\operatorname{span}}
\newcommand{\Ker}{\operatorname{Ker}}
\newcommand{\Ran}{\operatorname{Ran}}
\newcommand{\Cov}{\operatorname{Cov}}
\newcommand{\Var}{\operatorname{Var}}
\newcommand{\Corr}{\operatorname{Corr}}
\newcommand{\E}{\mathbb E}
\newcommand{\Prob}{\mathbb P}
\newcommand{\KL}{D_{\mathrm{KL}}}
\newcommand{\ESS}{\operatorname{ESS}}
\newcommand{\Rhat}{\widehat R}
\newcommand{\Id}{\operatorname{Id}}
\newcommand{\Hil}{\mathcal H}
\newcommand{\LF}{\ensuremath{\mathrm{LF}}}
\newcommand{\TMD}{\ensuremath{\mathrm{TMD}}}
\newcommand{\GTMD}{\ensuremath{\mathrm{GTMD}}}
\newcommand{\GPD}{\ensuremath{\mathrm{GPD}}}
\newcommand{\QCD}{\ensuremath{\mathrm{QCD}}}
\newcommand{\EFT}{\ensuremath{\mathrm{EFT}}}
\newcommand{\Amp}{\mathbf{Amp}}
\newcommand{\Dens}{\mathbf{Dens}}
\newcommand{\Match}{\mathbf{Match}}
\newcommand{\Red}{\mathbf{Red}}
\newcommand{\Proc}{\mathbf{Proc}}
\newcommand{\Infer}{\mathbf{Infer}}
\newcommand{\cO}{\mathcal O}
\newcommand{\cM}{\mathcal M}
\newcommand{\cR}{\mathcal R}
\newcommand{\cS}{\mathcal S}
\newcommand{\cV}{\mathcal V}
\newcommand{\cW}{\mathcal W}
\newcommand{\cE}{\mathcal E}
\newcommand{\cG}{\mathcal G}
\newcommand{\cK}{\mathcal K}
\newcommand{\cQ}{\mathcal Q}
\newcommand{\cP}{\mathcal P}
\newcommand{\cC}{\mathcal C}
\newcommand{\cZ}{\mathcal Z}
\newcommand{\cD}{\mathcal D}
\newcommand{\cN}{\mathcal N}
\newcommand{\cU}{\mathcal U}
\newcommand{\cF}{\mathcal F}
\newcommand{\cA}{\mathcal A}
\newcommand{\cB}{\mathcal B}
\newcommand{\cL}{\mathcal L}
\newcommand{\cJ}{\mathcal J}
\newcommand{\cI}{\mathcal I}
\newcommand{\cH}{\mathcal H}
\newcommand{\cT}{\mathcal T}
\newcommand{\cX}{\mathcal X}
\newcommand{\cY}{\mathcal Y}
\newcommand{\cRho}{\mathcal R}
\newcommand{\eps}{\varepsilon}
\newcommand{\bbone}{\bm 1}
\newcommand{\abs}[1]{\left|#1\right|}
\newcommand{\norm}[1]{\left\lVert#1\right\rVert}
\newcommand{\inner}[2]{\left\langle #1\middle|#2\right\rangle}
\newcommand{\avg}[1]{\left\langle #1\right\rangle}
\newcommand{\trans}{^{\mathsf T}}
\newcommand{\caldata}{\mathcal D_{\mathrm{cal}}}
\newcommand{\diagdata}{\mathcal D_{\mathrm{diag}}}
\newcommand{\holddata}{\mathcal D_{\mathrm{hold}}}
\newcommand{\allData}{\mathcal D}
\newcommand{\post}{\mathrm{post}}
\newcommand{\pred}{\mathrm{pred}}
\newcommand{\expdata}{\mathrm{exp}}
\newcommand{\latt}{\mathrm{latt}}
\newcommand{\thry}{\mathrm{th}}
\newcommand{\num}{\mathrm{num}}
\newcommand{\emul}{\mathrm{em}}
\newcommand{\disc}{\mathrm{disc}}
\newcommand{\sys}{\mathrm{sys}}
\newcommand{\stat}{\mathrm{stat}}
\newcommand{\hold}{\mathrm{hold}}
\newcommand{\val}{\mathrm{val}}
\newcommand{\obs}{\mathrm{obs}}
\newcommand{\repdata}{\mathrm{rep}}
\newcommand{\micro}{\mathrm{micro}}
\newcommand{\matchp}{\mathrm{match}}
\newcommand{\evol}{\mathrm{evol}}
\newcommand{\CS}{\ensuremath{\mathrm{CS}}}
\newcommand{\pert}{\mathrm{pert}}
\newcommand{\odd}{\mathrm{odd}}
\newcommand{\wil}{\mathrm{Wilson}}
\newcommand{\nuc}{\mathrm{nuc}}
\newcommand{\proc}{\mathrm{proc}}
\newcommand{\trunc}{\mathrm{trunc}}
\newcommand{\reg}{\mathrm{reg}}
\newcommand{\approxm}{\mathrm{approx}}
\newcommand{\score}{\mathcal S}
\newcommand{\GP}{\ensuremath{\mathrm{GP}}}
\newcommand{\HMC}{\ensuremath{\mathrm{HMC}}}
\newcommand{\NUTS}{\ensuremath{\mathrm{NUTS}}}
\newcommand{\SBC}{\ensuremath{\mathrm{SBC}}}
\newcommand{\LOO}{\ensuremath{\mathrm{LOO}}}
\newcommand{\PSIS}{\ensuremath{\mathrm{PSIS}}}

\newcolumntype{Y}{>{\raggedright\arraybackslash}X}
\newcolumntype{L}[1]{>{\raggedright\arraybackslash}p{#1}}

\title{\vspace{-1.0cm}
{\color{navy}\bfseries Volume VI: Shared Microscopic Inference,\\[2mm]
Correlated Uncertainty, and Withheld Validation}\\[4mm]
\large Likelihood construction, parameter ownership, model discrepancy,\\
posterior prediction, and falsifiable revision of the common light-front model}
\author{DeuteronWigner project}
\date{30 July 2026}

\begin{document}
\maketitle

\begin{center}
\begin{minipage}{0.94\textwidth}
\colorbox{softblue}{\parbox{0.96\textwidth}{
\textbf{Status and purpose.}
Volumes~I--V define the regulated microscopic state spaces, common nucleon
overlaps, Wilson-line phases, matched spin-1 nuclear operators, and the
regulator-to-QCD evolution and process layer.  The present volume defines how
those objects may be calibrated without rebuilding the earlier phenomenological
patchwork in statistical form.  The inferential state is one shared posterior
over Hamiltonian, current, Wilson, nuclear, matching, evolution, nuisance, and
controlled-discrepancy parameters.  Every posterior member retains a common
microscopic identity through every flavor, target polarization, parton
species, reduction, scale, nuclear mechanism, and process observable.
Calibration and withheld validation are fixed before final optimization.
Failure of a withheld observable must revise a physical interaction, state
sector, current, matching map, factorization assumption, or declared
discrepancy model; it may not be repaired by adding an isolated coefficient to
the failed TMD.
}}
\end{minipage}
\end{center}

\begin{abstract}
A generative light-front model becomes scientifically predictive only when a
restricted set of shared physical parameters is confronted with heterogeneous
data and then succeeds on observables that were not used to determine those
parameters.  This volume constructs the inference, uncertainty, and
falsifiability layer for the DeuteronWigner program.  The central object is a
typed probabilistic program whose deterministic forward map is the complete
chain from regulated Hamiltonian and current parameters through common GTMD
operators, Wilson-line dynamics, matched nuclear composition, QCD matching,
rank-aware evolution, and process-qualified observables.

The parameter vector is divided by physical ownership: Hamiltonian vertices,
counterterms, currents, Wilson kernels, nuclear interactions, operator
matching, Collins--Soper and perturbative inputs, experimental nuisance
parameters, and explicit discrepancy operators.  A parameter cannot migrate
into one independent normalization or width per named TMD.  Exact symmetries,
support, positivity, conservation laws, and composition exclusions are
enforced structurally; priors are used for uncertain physical scales and
naturalness, not to imitate exact constraints.

We formulate likelihoods for event counts, binned cross sections and
asymmetries, correlated lattice matrix elements, form factors, extracted
partonic information, and theory-calibration observables.  Dataset ancestry is
part of identity so raw data and quantities extracted from the same events are
not counted twice.  Correlated experimental systematics are represented by
covariance matrices or explicit nuisance parameters.  Theory uncertainty is
decomposed into Hamiltonian EFT truncation, Fock/basis truncation,
regulator/renormalization, Wilson order, LF-to-QCD matching, TMD evolution,
nuclear many-body, process-power, emulator, and numerical components.
Correlated truncation fields are tied to power counting and operator ownership;
they are not an unrestricted catch-all Gaussian process.

The posterior is tested through prior predictive simulation, synthetic
parameter recovery, simulation-based calibration, exact-versus-emulated
comparisons, multi-chain diagnostics, posterior predictive checks,
leave-family-out prediction, proper scoring rules, and predefined withheld
observable gates.  The volume distinguishes calibration fit, model criticism,
and scientific validation.  It defines strong cross-channel tests, such as
using a common rescattering sector constrained by selected Sivers information
to predict Boer--Mulders, tensor naive-\(T\)-odd, antiquark, or gluon
color-sector observables without separate retuning.  A final section turns the
posterior into an experimental-design object by maximizing expected information
gain or a controlled Fisher approximation while retaining feasibility and
systematic constraints.

The result is a fail-closed inference architecture in which uncertainty bands
are pushforwards of common microscopic members, model discrepancy remains
physically attributable, and a failed withheld family produces a traceable
model-revision decision rather than a new curve-specific parameter.
\end{abstract}

\tableofcontents

\section{Scope, interfaces, and status language}
\label{sec:scope}

\subsection{Input contract}

The inference layer receives a deterministic or conditionally deterministic
forward model
\begin{equation}
 \cG:
 (\theta,\nu,\eta_{\sys})
 \longmapsto
 \left\{
 \cO_i^{\thry}
 \right\}_{i\in\cI_{\obs}},
 \label{eq:forward-map}
\end{equation}
where:
\begin{itemize}
\item \(\theta\) denotes shared continuous physical and matching parameters;
\item \(\nu\) denotes discrete, typed model choices such as regulator family,
      Fock truncation, nuclear interaction, Wilson order, or perturbative
      accuracy;
\item \(\eta_{\sys}\) denotes experimental or external-input nuisance
      parameters;
\item each \(\cO_i^{\thry}\) carries complete operator, scale, process,
      provenance, and numerical-accuracy metadata.
\end{itemize}
The map in \cref{eq:forward-map} is not a flat table of unrelated predictions.
It is the composed chain
\begin{equation}
 \begin{tikzcd}[column sep=small]
 \{\theta_H,\theta_J,\theta_{\wil}\}
 \arrow[r]
 & \{|\Psi_N\rangle,|\Psi_D\rangle,\cO_{\GTMD}\}
 \arrow[r]
 & \{\theta_{\matchp},\theta_{\evol}\}
 \arrow[r]
 & \{\widetilde F_A(\mu,\zeta),\cO_{\proc}\}
 \arrow[r]
 & \{\cO_i^{\thry}\}.
 \end{tikzcd}
 \label{eq:generative-inference-chain}
\end{equation}
Every intermediate object is versioned.  A posterior member must be replayable
through the same composition graph.

The data are divided into
\begin{equation}
 \allData
 =
 \caldata
 \,\dot\cup\,
 \diagdata
 \,\dot\cup\,
 \holddata,
 \label{eq:data-partition}
\end{equation}
where \(\caldata\) determines the posterior, \(\diagdata\) supports model and
computational criticism without being used as a final predictive claim, and
\(\holddata\) is protected for scientific validation.

\begin{requirement}[Frozen split, V6.DATA.1]
The membership and ancestry of \(\caldata\), \(\diagdata\), and \(\holddata\)
must be committed before final calibration.  Moving a failed withheld family
into the calibration set changes the scientific question and requires a new
versioned analysis.
\end{requirement}

\subsection{Output contract}

The principal inferential output is
\begin{equation}
 p(\Theta,\nu,\eta,\delta\mid\caldata,\mathfrak M),
 \qquad
 \Theta=
 \left(
 \theta_H,\theta_J,\theta_{\wil},\theta_{\nuc},
 \theta_{\matchp},\theta_{\evol}
 \right),
 \label{eq:full-posterior}
\end{equation}
where \(\mathfrak M\) is the complete model and convention manifest and
\(\delta\) denotes controlled discrepancy fields.

The output is stored as a common member ensemble
\begin{equation}
 \mathcal E_{\post}
 =
 \left\{
 e_s=
 \left(
 \Theta_s,\nu_s,\eta_s,\delta_s,
 \mathrm{state\ hashes},
 \mathrm{operator\ hashes},
 \mathrm{matching\ manifest},
 \mathrm{process\ manifest}
 \right)
 \right\}_{s=1}^{N_{\post}}.
 \label{eq:posterior-ensemble}
\end{equation}
A prediction for any observable family \(\cY_*\) is the pushforward
\begin{equation}
 p(\cY_*\mid\caldata)
 =
 \int
 p(\cY_*\mid\Theta,\nu,\eta,\delta)
 p(\Theta,\nu,\eta,\delta\mid\caldata)
 \,\dd\Theta\,\dd\nu\,\dd\eta\,\dd\delta .
 \label{eq:posterior-predictive}
\end{equation}

\begin{requirement}[Shared-member output, V6.ENS.1]
Independent marginal bands generated with unrelated random seeds or parameter
draws are not an admissible substitute for \(\mathcal E_{\post}\).  Cross-TMD,
cross-flavor, quark--gluon, proton--neutron, nucleon--deuteron, scale, and
process covariances must be recoverable from shared member identity.
\end{requirement}

\subsection{What this volume does not claim}

This volume does not claim:
\begin{itemize}
\item that a Bayesian posterior converts an incomplete Hamiltonian into QCD;
\item that all uncertainties are subjective or interchangeable;
\item that a good fit to \(\caldata\) validates the microscopic model;
\item that an unrestricted discrepancy field may absorb missing operators;
\item that a credible interval has frequentist coverage without a coverage
      study;
\item that a Gaussian likelihood is appropriate for every dataset;
\item that a surrogate is trustworthy outside its validated domain;
\item that Bayes factors or one scalar information criterion decide scientific
      adequacy;
\item that model averaging can make mutually inconsistent descriptions
      simultaneously true;
\item that a withheld prediction is genuine when its normalization or shape
      was indirectly tuned using the same underlying events.
\end{itemize}

\subsection{Evidence and uncertainty classes}

The inferential layer retains the evidence vocabulary of the preceding
volumes and adds a statistical-status field:
\begin{longtable}{L{0.22\textwidth}L{0.69\textwidth}}
\toprule
Status & Meaning\\
\midrule
\endhead
\texttt{EXACT} & Algebraic, symmetry, kinematic, support, conservation, or
typed-composition consequence.  It is not assigned a prior width.\\
\texttt{CALIBRATION\_DATA} & Observation or matrix element included in
\(\caldata\).\\
\texttt{DIAGNOSTIC\_DATA} & Observation reserved for criticism or workflow
development but not advertised as a withheld prediction.\\
\texttt{WITHHELD\_DATA} & Predeclared family excluded from posterior
construction and used for scientific prediction.\\
\texttt{FIT/PHENOMENOLOGY} & External extraction or fit whose covariance and
scheme are available and whose data ancestry is declared.\\
\texttt{LATTICE\_INFORMED} & Lattice matrix element with stated action,
spacing, volume, matching, and covariance provenance.\\
\texttt{MODEL} & Replaceable physical approximation with a named parameter
and validation path.\\
\texttt{DISCREPANCY} & Explicit residual operator or stochastic field with a
declared physical scope, power counting, and correlation model.\\
\texttt{NUMERICAL} & Eigensolver, quadrature, transform, interpolation,
emulator, or Monte Carlo error.\\
\texttt{UNASSESSED} & Required information is absent.  The object cannot enter
an authoritative posterior or prediction.\\
\bottomrule
\end{longtable}

\begin{principle}[Inference cannot upgrade evidence]
Posterior concentration does not convert a model assumption into a fitted
operator identity, and a structural acceptance gate does not convert a model
prediction into a measurement.  Evidence class and inferential status are
stored separately.
\end{principle}

\section{The typed probabilistic program}
\label{sec:program}

\subsection{Parameter blocks and physical ownership}

The complete continuous parameter vector is partitioned by physical owner:
\begin{equation}
 \Theta=
 \left(
 \theta_H,\theta_{\mathrm{ct}},
 \theta_J,\theta_{\wil},
 \theta_{\nuc},\theta_{\matchp},
 \theta_{\CS},\theta_{\proc}
 \right).
 \label{eq:parameter-blocks}
\end{equation}
The blocks mean:
\begin{longtable}{L{0.18\textwidth}L{0.71\textwidth}}
\toprule
Block & Physical ownership\\
\midrule
\endhead
\(\theta_H\) & Masses and interaction coefficients in the regulated
light-front Hamiltonian.\\
\(\theta_{\mathrm{ct}}\) & Counterterms, sector-dependent renormalization
coefficients, and truncation-induced operators.\\
\(\theta_J\) & One-body, exchange, induced, and local currents consistent
with the Hamiltonian.\\
\(\theta_{\wil}\) & Eikonal vertices, infrared scales, Wilson-order
coefficients, and shared rescattering interactions.\\
\(\theta_{\nuc}\) & Nuclear interaction and operator coefficients, including
explicit higher-sector and coherent amplitudes.\\
\(\theta_{\matchp}\) & LF-to-QCD, small-\(b\), threshold, and operator-mixing
matching coefficients not fixed perturbatively.\\
\(\theta_{\CS}\) & Nonperturbative Collins--Soper and evolution-boundary
parameters shared by the appropriate color representation.\\
\(\theta_{\proc}\) & Process-specific quantities only when they are not fixed
by a theorem or independent perturbative calculation.\\
\bottomrule
\end{longtable}

Each parameter has an ownership record
\begin{equation}
 \operatorname{Owner}(\theta_a)
 =
 \left(
 \begin{gathered}
 \text{physical object},\ \text{operator block},\ \text{symmetry class},\\
 \text{sector support},\ \text{scale dependence},\
 \text{observables reached}
 \end{gathered}
 \right).
 \label{eq:owner-record}
\end{equation}

\begin{requirement}[No parameter migration, V6.PARAM.1]
A parameter may enter an observable only through the physical object that owns
it.  A failed tensor TMD cannot acquire a new coefficient unless that
coefficient is attached to a missing Hamiltonian interaction, state sector,
current, Wilson kernel, matching operator, or declared discrepancy object that
also propagates to every other observable reached by the same physics.
\end{requirement}

\subsection{Deterministic graph and stochastic nodes}

Let \(G=(V,E)\) be the typed composition graph.  Each deterministic node
\(v\in V_{\det}\) computes
\begin{equation}
 x_v=f_v\left(\{x_u:u\to v\},\Theta_v,\nu_v\right),
 \label{eq:det-node}
\end{equation}
while stochastic nodes include priors, nuisance variables, discrepancy fields,
emulator residuals, and observations.  The joint density factorizes as
\begin{equation}
 p(\allData,\Theta,\nu,\eta,\delta)
 =
 p(\Theta,\nu)\,
 p(\eta)\,
 p(\delta\mid\Theta,\nu)
 \prod_{d\in\allData}
 p\!\left(
 y_d\mid
 \cG_d(\Theta,\nu),\eta_d,\delta_d
 \right).
 \label{eq:joint-factorization}
\end{equation}

The graph contains explicit exclusion relations.  Mutually exclusive
descriptions, such as an explicit \(qqqg\) sector and the induced operator that
replaces it after elimination, are not independent stochastic summands.

\begin{proposition}[Posterior provenance]
If every stochastic parameter has one owner, every deterministic map is
versioned, and every data node records its ancestry, then every posterior
correlation between two predictions has a traceable path through shared
physical or nuisance variables.
\end{proposition}

\begin{proof}
Both predictions are measurable functions of the common stochastic nodes in
\cref{eq:joint-factorization}.  Their covariance is generated by nodes in the
intersection of their ancestor sets.  Versioned graph traversal identifies
that intersection and therefore the origin of the covariance.
\end{proof}

\subsection{Discrete model and truncation labels}

The discrete object
\begin{equation}
 \nu=
 \left(
 r_H,\mathcal F_H,N_{\max},
 n_{\wil},r_{\nuc},
 \mathfrak s_{\TMD},\mathfrak A_{\pert},
 \mathfrak P_{\proc}
 \right)
 \label{eq:discrete-label}
\end{equation}
contains regulator, Fock, basis, Wilson-order, nuclear-resolution, TMD-scheme,
perturbative-accuracy, and process records.

Discrete alternatives are classified as:
\begin{itemize}
\item \emph{nested truncations}, with a declared convergence direction;
\item \emph{equivalent representations}, related by matching or a unitary map;
\item \emph{competing physical hypotheses}, such as different missing
      interactions;
\item \emph{external model variants}, useful for sensitivity but not assumed
      to be draws from one probability distribution.
\end{itemize}

A wave-function envelope is not automatically a posterior over wave functions.
A discrete prior is scientifically meaningful only when its support and
weights have a stated interpretation.

\subsection{Exact constraints are not narrow priors}

Support, color singletness, Hermiticity, current conservation, target-spin
closure, positivity where valid, and no-double-counting relations are built
into the parameterization or graph.  For example, if
\begin{equation}
 \Phi=LL^\dagger ,
 \label{eq:structural-positive}
\end{equation}
then positivity is exact for every parameter value.  It should not be imposed
by a narrow penalty on negative eigenvalues.

\begin{requirement}[Structural constraints, V6.PARAM.2]
An exact relation must be enforced by coordinates, projectors, constrained
maps, or deterministic transformations.  Priors and likelihood penalties are
reserved for uncertain statements.
\end{requirement}

\section{Data identity, ancestry, and protected partitions}
\label{sec:data}

\subsection{Dataset object}

A dataset is not identified by a filename alone.  The minimal record is
\begin{equation}
 \mathfrak D_d=
 \left(
 \begin{gathered}
 \text{observable definition},\ \text{process},\ \text{kinematics},\
 \text{units},\ \text{binning},\\
 \text{polarization convention},\ \text{scheme},\ \text{central values},\
 \text{covariance},\\
 \text{nuisance map},\ \text{event ancestry},\ \text{selection},\
 \text{partition}
 \end{gathered}
 \right).
 \label{eq:dataset-object}
\end{equation}
The event-ancestry field identifies whether two published quantities are
derived from the same events, unfolding, normalization sample, lattice gauge
ensemble, PDF fit, or common external calibration.

\begin{requirement}[No data double counting, V6.DATA.2]
Two datasets with overlapping ancestry may enter one likelihood only through a
joint covariance or a hierarchical source model.  Treating a raw cross
section, an asymmetry extracted from it, and a PDF fit containing the same
measurement as independent observations is forbidden.
\end{requirement}

\subsection{Families rather than random points}

The partition in \cref{eq:data-partition} is made at the level of scientifically
coherent families.  Examples include:
\begin{itemize}
\item an entire experiment, target polarization, or beam-energy setting;
\item one observable class such as tensor DIS, SIDIS Sivers, or tagged
      spectator data;
\item a group of correlated lattice moments from one ensemble;
\item one flavor or parton channel;
\item a kinematic region that tests extrapolation rather than interpolation;
\item a process with a different Wilson-link or color-flow class.
\end{itemize}

Randomly withholding individual points from a densely correlated curve tests
interpolation and may substantially overstate predictive ability.

\begin{definition}[Calibration, diagnostic, and withheld roles]
\(\caldata\) is permitted to alter the posterior.  \(\diagdata\) may be used
during model development to reveal defects, but any subsequent claim using it
is labeled diagnostic rather than withheld.  \(\holddata\) remains inaccessible
to parameter tuning and model selection until the analysis version is frozen.
\end{definition}

\subsection{Candidate partition for the first microscopic analysis}

A credible initial partition is:
\begin{longtable}{L{0.21\textwidth}L{0.34\textwidth}L{0.35\textwidth}}
\toprule
Layer & Calibration candidates & Withheld candidates\\
\midrule
\endhead
Spectrum/state & Selected hadron masses, charges, magnetic and axial moments &
Unused current components, higher states, rotational multiplets.\\
Nucleon collinear & Restricted unpolarized/helicity moments and selected
flavor information & Flavor-separated curves or moments outside the
calibration region.\\
Off-forward & Selected electromagnetic form factors & Unused GPD or EMT
moments and OAM-sensitive relations.\\
Wilson sector & A restricted Sivers subset or one process sign & Boer--Mulders,
antiquark, tensor-\(T\)-odd, and gluon \(f/d\) observables.\\
Deuteron & Binding, asymptotic S/D information, static moments, selected
elastic points & \(b_1\), tagged tensor observables, unused elastic
combinations, non-nucleonic-sensitive channels.\\
Evolution/process & A restricted multi-\(Q\) family in an established channel &
Another scale, process, rank, or color-flow family.\\
\bottomrule
\end{longtable}
The actual split is committed in a machine-readable manifest and may differ
from this example.

\subsection{Blinding and access control}

For a strong withheld test, the software should support:
\begin{itemize}
\item encrypted or separately stored withheld values;
\item metadata access without numerical values, sufficient to construct
      predictions before unblinding;
\item a hash of the split and validation metric definitions;
\item an immutable timestamp and analysis commit;
\item a rule for handling unexpected data-quality problems without silently
      redefining the target.
\end{itemize}

\section{Likelihood construction}
\label{sec:likelihood}

\subsection{General observation model}

For dataset \(d\), let \(\bm y_d\) be the measured vector and
\(\bm m_d(\Theta,\nu)\) the model prediction after bin integration, acceptance,
radiative treatment, and declared process assembly.  A general observation
model is
\begin{equation}
 \bm y_d
 =
 \bm m_d(\Theta,\nu)
 +
 \bm A_d\eta_d
 +
 \bm\delta_d
 +
 \bm\epsilon_d,
 \label{eq:observation-model}
\end{equation}
where \(\bm A_d\eta_d\) represents explicit correlated nuisance shifts,
\(\bm\delta_d\) is controlled theory discrepancy, and
\(\bm\epsilon_d\) is measurement noise.

The likelihood is defined at the least processed level for which the full
measurement model is available.  A model prediction at bin center is not
compared with a bin average unless the corresponding approximation is
quantified.

\subsection{Correlated Gaussian measurements}

For approximately Gaussian data,
\begin{equation}
 p(\bm y_d\mid\Theta,\eta_d,\delta_d,\nu)
 =
 \mathcal N\!\left(
 \bm y_d\ ;\
 \bm m_d(\Theta,\nu)+\bm A_d\eta_d+\bm\delta_d,\,
 \bm C_{d,\stat}
 \right),
 \label{eq:gaussian-likelihood}
\end{equation}
with nuisance prior
\begin{equation}
 \eta_d\sim\mathcal N(\bm 0,\bm C_{d,\sys}).
 \label{eq:nuisance-prior}
\end{equation}
Marginalizing \(\eta_d\) gives covariance
\begin{equation}
 \bm C_d
 =
 \bm C_{d,\stat}
 +
 \bm A_d\bm C_{d,\sys}\bm A_d\trans ,
 \label{eq:marginal-covariance}
\end{equation}
when the nuisance response is linear and Gaussian.

\begin{proposition}[Nuisance/covariance equivalence]
Under \cref{eq:gaussian-likelihood,eq:nuisance-prior}, explicit Gaussian
nuisance marginalization and the covariance form
\cref{eq:marginal-covariance} produce the same likelihood.  The explicit
nuisance form remains preferable when normalization, acceptance, or
polarization responses are nonlinear or shared across datasets.
\end{proposition}

\subsection{Multiplicative normalization and polarization}

For a normalization nuisance \(n_d\),
\begin{equation}
 \bm m_d^{\sys}
 =
 \left(1+n_d\right)\bm m_d,
 \qquad
 n_d\sim\mathcal N(0,\sigma_{n_d}^2),
 \label{eq:multiplicative-normalization}
\end{equation}
or a log-normal parameterization is used when positivity and large relative
errors matter.  Beam and target polarizations, dilution factors, luminosity,
and fragmentation normalizations remain separate nuisance parameters when
their correlations differ.

The model must distinguish an experimental normalization uncertainty from a
physical normalization parameter.  A nuisance cannot be absorbed into a TMD
normalization and then propagated as microscopic information.

\subsection{Count and event likelihoods}

For independent event counts,
\begin{equation}
 N_b\sim\operatorname{Poisson}\!\left(\lambda_b(\Theta,\eta,\nu)\right),
 \label{eq:poisson-likelihood}
\end{equation}
with
\begin{equation}
 \lambda_b
 =
 \mathcal L
 \int_{\mathrm{bin}\ b}
 \dd\Phi\,
 \epsilon(\Phi,\eta)\,
 \frac{\dd\sigma(\Theta,\nu)}{\dd\Phi}
 +\lambda_{b,\mathrm{bkg}} .
 \label{eq:expected-count}
\end{equation}
For spin-state counts, the joint likelihood retains the actual polarization,
luminosity, live-time, and background structure.  A Gaussian asymmetry
likelihood is used only when its approximation and covariance are adequate.

Unbinned event likelihoods are allowed when acceptance and normalization are
well defined:
\begin{equation}
 \log\cL
 =
 -\Lambda(\Theta)
 +\sum_{e=1}^{N}
 \log\lambda(\Phi_e\mid\Theta).
 \label{eq:extended-unbinned}
\end{equation}

\subsection{Lattice and Euclidean matrix elements}

A lattice dataset records the gauge ensembles, lattice spacing, volume, quark
masses, renormalization, continuum/chiral extrapolation, and covariance.
Correlated lattice observables enter as
\begin{equation}
 \bm y_{\latt}
 \sim
 \mathcal N\!\left(
 \bm Z_{\latt\to\mathfrak s}\,
 \bm m_{\latt}(\Theta)
 +\bm\delta_{\latt},
 \bm C_{\latt}
 \right).
 \label{eq:lattice-likelihood}
\end{equation}
Matching or extrapolation uncertainty is not added twice when it is already
contained in \(\bm C_{\latt}\) or represented by explicit nuisance variables.

Published quasi- or pseudo-distributions are used as calibration data only
when their matching, covariance, and common-ensemble correlations are
available.  Otherwise they are diagnostic constraints or moments.

\subsection{PDF and TMD extractions as data}

An external PDF/TMD ensemble is not a set of independent measurements.  Its
use requires:
\begin{itemize}
\item the operator and factorization scheme;
\item the scale and perturbative order;
\item the data ancestry of the extraction;
\item a member-level covariance or a reproducible likelihood interface;
\item separation from any raw data also used directly.
\end{itemize}
When these are unavailable, the extraction is treated as a comparison
ensemble, not as a Gaussian likelihood built from pointwise error bars.

\subsection{Asymmetric and non-Gaussian uncertainty}

Split-normal, Student-\(t\), skew-normal, mixture, or direct likelihoods may be
used when justified.  The choice is part of dataset identity.  Replacing a
published asymmetric interval by a symmetric standard deviation is a model
assumption and must be validated.

\subsection{Likelihood normalization, units, and transformations}

Every likelihood reports:
\begin{itemize}
\item units and Jacobians under variable transformations;
\item whether constants independent of parameters are retained;
\item bin widths and phase-space factors;
\item covariance regularization, if any;
\item masked bins and selection criteria;
\item the numerical log-likelihood tolerance.
\end{itemize}

\begin{requirement}[Likelihood audit, V6.LIKE.1]
The same synthetic dataset evaluated through two algebraically equivalent
likelihood paths must agree within numerical tolerance.  Unit, bin, sign, and
normalization mismatches are injected as mandatory failures.
\end{requirement}

\section{Prior hierarchy and physical regularization}
\label{sec:priors}

\subsection{Prior factorization}

The prior is organized by ownership:
\begin{align}
 p(\Theta,\nu,\eta,\delta)
 ={}&
 p(\nu)\,
 p(\theta_H,\theta_{\mathrm{ct}}\mid\nu)\,
 p(\theta_J\mid\theta_H,\nu)\,
 p(\theta_{\wil}\mid\theta_H,\nu)
 \nonumber\\
 &\times
 p(\theta_{\nuc}\mid\theta_H,\nu)\,
 p(\theta_{\matchp},\theta_{\CS}\mid\nu)\,
 p(\eta)\,
 p(\delta\mid\Theta,\nu).
 \label{eq:prior-factorization}
\end{align}
The factorization is not assumed to imply physical independence.  Shared
hyperparameters encode common naturalness, symmetry breaking, regulator
running, or representation dependence.

\subsection{Naturalness and power counting}

For a dimensionless EFT or effective-Hamiltonian coefficient \(c_a\),
\begin{equation}
 c_a\mid\bar c,\rho
 \sim
 \mathcal N(0,\bar c^2 \rho_a^2),
 \label{eq:naturalness-prior}
\end{equation}
where \(\rho_a\) records known symmetry or operator scaling and \(\bar c\) may
have a weakly informative hyperprior.  Coefficients known to be nonzero at a
lower order are centered on that lower-order information rather than zero.

Naturalness priors reduce overfitting only when the operator basis and
dimensionless scaling are physically declared.  They do not license arbitrary
shrinkage of all parameters toward one common magnitude
\cite{Wesolowski2016,Furnstahl2015}.

\subsection{Regulator-dependent coefficients}

Counterterms at regulator \(r\) are not exchangeable draws across cutoffs.
Their prior is conditional on the renormalization trajectory:
\begin{equation}
 \theta_{\mathrm{ct}}^{(r')}
 =
 \cR_{r'\leftarrow r}
 \theta_{\mathrm{ct}}^{(r)}
 +
 \xi_{r'r},
 \label{eq:counterterm-flow-prior}
\end{equation}
where \(\xi_{r'r}\) represents residual truncation after the declared matching
map.  A prior that forces bare coefficients to remain numerically identical
under regulator changes is generally inappropriate.

\subsection{Symmetry-breaking hierarchy}

Charge-symmetry, rotational, chiral, heavy-flavor, or tensor-breaking
parameters carry power-counting hyperparameters.  For example,
\begin{equation}
 \theta_{\mathrm{CSB},a}
 =
 \epsilon_{\mathrm{CSB}}\,
 \widetilde\theta_a,
 \qquad
 \widetilde\theta_a\sim\mathcal N(0,\bar c_{\mathrm{CSB}}^2),
 \label{eq:csb-hierarchy}
\end{equation}
with \(\epsilon_{\mathrm{CSB}}\) tied to the declared physical expansion.  The
same hyperparameter propagates to all operators in the symmetry class.

\subsection{Prior predictive checks}

Before conditioning on \(\caldata\), draw
\begin{equation}
 \Theta^{(s)},\nu^{(s)},\eta^{(s)},\delta^{(s)}
 \sim p(\Theta,\nu,\eta,\delta),
 \qquad
 \bm y_{\repdata}^{(s)}
 \sim p(\bm y\mid\Theta^{(s)},\nu^{(s)},\eta^{(s)},\delta^{(s)}).
 \label{eq:prior-predictive}
\end{equation}
Prior predictive checks must expose:
\begin{itemize}
\item unphysical spectra or negative rates;
\item implausible charges, moments, and form factors;
\item excessive or negligible symmetry breaking;
\item Wilson phases too large for the declared order;
\item matching coefficients that violate known limits;
\item discrepancy fields broader than the physics they are meant to
      approximate;
\item numerical regions outside the solver or emulator domain.
\end{itemize}

\begin{requirement}[Prior predictive admissibility, V6.PRIOR.1]
A prior is not accepted merely because the posterior is numerically stable.
It must generate scientifically plausible observables before seeing
\(\caldata\) and must not place substantial probability on states rejected by
exact structural gates.
\end{requirement}

\subsection{Weakly informative does not mean unowned}

A weakly informative prior still records:
\begin{equation}
 \left(
 \text{units},
 \text{reference scale},
 \text{operator owner},
 \text{support},
 \text{tail behavior},
 \text{hyperparameters},
 \text{sensitivity plan}
 \right).
 \label{eq:prior-record}
\end{equation}
Changing a prior after inspecting \(\holddata\) invalidates the withheld claim.
Prior sensitivity is assessed on calibration and diagnostic quantities or in
a new versioned analysis.

\section{Theory discrepancy and correlated truncation}
\label{sec:discrepancy}

\subsection{Why discrepancy is necessary but dangerous}

The observation model
\begin{equation}
 y(x)=\cG(x;\Theta)+\delta(x)+\epsilon(x)
 \label{eq:basic-discrepancy}
\end{equation}
separates simulator output, model discrepancy, and observation error.
Ignoring \(\delta\) can produce biased and overconfident physical parameters,
while an unrestricted \(\delta\) can become confounded with those parameters
\cite{KennedyOHagan2001,BrynjarsdottirOHagan2014,Plumlee2017}.

For this program, discrepancy is not one universal function.  It is a direct
sum of physically typed components:
\begin{equation}
 \delta
 =
 \delta_H
 +\delta_{\mathrm{Fock}}
 +\delta_{\reg}
 +\delta_{\wil}
 +\delta_{\matchp}
 +\delta_{\evol}
 +\delta_{\nuc}
 +\delta_{\proc}
 +\delta_{\emul}
 +\delta_{\num}.
 \label{eq:discrepancy-decomposition}
\end{equation}
The components may be correlated, but their meanings remain distinct.

\subsection{Hamiltonian and EFT truncation}

For an observable with expansion
\begin{equation}
 \cO^{(k)}(x)
 =
 \cO_{\mathrm{ref}}(x)
 \sum_{n=0}^{k} c_n(x)Q(x)^n,
 \label{eq:eft-expansion}
\end{equation}
the truncation remainder is
\begin{equation}
 \delta_H^{(k)}(x)
 =
 \cO_{\mathrm{ref}}(x)
 \sum_{n=k+1}^{\infty} c_n(x)Q(x)^n.
 \label{eq:eft-remainder}
\end{equation}
A correlated model places a process on the coefficient functions:
\begin{equation}
 c_n(x)\sim
 \GP\!\left(
 \mu_n(x),
 \bar c^2 K_n(x,x';\ell)
 \right),
 \label{eq:eft-gp}
\end{equation}
with correlations across kinematics and, when justified, across observables
that share the same operator expansion.  This implements correlated EFT
truncation rather than independent pointwise error bars
\cite{Furnstahl2015,Melendez2019}.

The expansion parameter, reference scale, first omitted order, and correlation
length are inferred or varied only within the EFT interpretation.  A GP
kernel is not chosen solely for visual smoothness.

\subsection{Fock-sector and basis truncation}

Let \(\cO_{r,n}\) denote the matched observable at regulator/basis/Fock level
\((r,n)\).  The residual
\begin{equation}
 \Delta_{r,n}
 =
 \cO_{r,n}
 -
 \cR_{r,n\leftarrow r',n'}[\cO_{r',n'}]
 \label{eq:fock-basis-residual}
\end{equation}
is modeled against known convergence variables such as \(1/N_{\max}\),
longitudinal resolution, omitted-sector probability scale, or ultraviolet
cutoff.  Bare wave-function differences are not substituted for observable
discrepancy.

\subsection{Wilson-order discrepancy}

If
\begin{equation}
 W_{\odd}^{(N)}
 =
 \sum_{n=1}^{N}W_{\odd}^{[n]},
 \label{eq:wilson-order-series}
\end{equation}
then
\begin{equation}
 \delta_{\wil}^{(N)}
 =
 W_{\odd}-W_{\odd}^{(N)}
 \label{eq:wilson-remainder}
\end{equation}
is correlated among every Sivers-like, Boer--Mulders-like, tensor, antiquark,
and gluon projection that uses the same omitted Wilson webs.  It cannot be
represented by independent phase errors for each TMD.

\subsection{Matching and evolution discrepancy}

The LF-to-QCD matching residual is an operator-space field,
\begin{equation}
 \widetilde{\bm W}^{\,\QCD}
 =
 \bm Z_{\LF\to\QCD}
 \otimes_x
 \widetilde{\bm W}^{\,\LF}
 +
 \bm\delta_{\matchp},
 \label{eq:matching-discrepancy}
\end{equation}
whose covariance respects operator mixing, rank, link/color class, and local
limits.

Evolution discrepancy is separated into:
\begin{equation}
 \delta_{\evol}
 =
 \delta_{\mathrm{pert\ order}}
 +
 \delta_{\mathrm{finite\ curl}}
 +
 \delta_{\mathrm{large}\ b}
 +
 \delta_{\mathrm{threshold}}
 +
 \delta_{W+Y}.
 \label{eq:evolution-discrepancy}
\end{equation}
Scale variation is a diagnostic or prior input for these terms; it is not
automatically a calibrated probability distribution.

\subsection{Nuclear and process discrepancy}

Nuclear discrepancy is attached to omitted many-body operators, sector
truncation, coherent amplitudes, or scale separation.  It retains target
helicity and operator structure.  A scalar unpolarized nuclear ratio cannot
serve as tensor or polarized discrepancy unless a theorem supplies that
relation.

Process discrepancy contains declared power corrections, fragmentation
uncertainty, missing fixed-order terms, and factorization limitations.
A channel marked \texttt{BROKEN} or \texttt{UNASSESSED} does not acquire a
Gaussian discrepancy that magically restores standard TMD factorization.

\subsection{Orthogonality and identifiability}

Let
\begin{equation}
 J_a(x)=\frac{\partial\cG(x;\Theta)}{\partial\theta_a}.
 \label{eq:model-tangent}
\end{equation}
A discrepancy basis \(\{B_m(x)\}\) can be constrained to satisfy
\begin{equation}
 \langle B_m,J_a\rangle_{\bm C^{-1}}=0
 \label{eq:orthogonal-discrepancy}
\end{equation}
for selected physically interpretable parameters.  This reduces confounding
between calibration directions and discrepancy, while not implying that the
true discrepancy is exactly orthogonal
\cite{Plumlee2017}.

Alternative safeguards include:
\begin{itemize}
\item informative priors from power counting;
\item multiple observable families with different parameter sensitivities;
\item explicit local-current and sum-rule constraints;
\item discrepancy support restricted to known omitted sectors or kinematic
      domains;
\item sensitivity analyses under several plausible discrepancy models.
\end{itemize}

\begin{requirement}[No catch-all repair, V6.DISC.1]
A discrepancy component must declare its physical source, operator support,
correlation structure, expected order, and observables reached.  A
curve-specific unconstrained residual added after a failed withheld test is not
an admissible model update.
\end{requirement}

\subsection{Law of total covariance}

For predictions \(\bm O\), the total covariance is
\begin{align}
 \Cov(\bm O)
 ={}&
 \E_{\nu,\delta}\!\left[
 \Cov(\bm O\mid\nu,\delta,\caldata)
 \right]
 \nonumber\\
 &+
 \Cov_{\nu,\delta}\!\left(
 \E[\bm O\mid\nu,\delta,\caldata]
 \right).
 \label{eq:total-covariance}
\end{align}
Within a fixed model,
\begin{equation}
 \Cov(\bm O\mid\nu)
 =
 \Cov_{\Theta,\eta}(\bm O)
 +
 \Cov_{\disc}(\bm O)
 +
 \Cov_{\emul}(\bm O)
 +
 \Cov_{\num}(\bm O)
 +
 \text{cross terms},
 \label{eq:covariance-components}
\end{equation}
where cross terms are retained whenever the components are not conditionally
independent.

\begin{principle}[Separated reporting]
The total posterior predictive distribution is authoritative for prediction,
but the Hamiltonian, Fock, regulator, Wilson, matching, evolution, nuclear,
process, experimental, emulator, and numerical components must remain
separately queryable.  One category cannot hide nonconvergence in another.
\end{principle}

\section{Posterior construction and hierarchical calibration}
\label{sec:posterior}

\subsection{Full posterior}

The posterior for one declared model version is
\begin{align}
 &p(\Theta,\nu,\eta,\delta,\lambda\mid\caldata)
 \nonumber\\
 &\quad\propto
 p(\caldata\mid\Theta,\nu,\eta,\delta)
 p(\Theta\mid\nu,\lambda)
 p(\delta\mid\Theta,\nu,\lambda)
 p(\eta)
 p(\nu)
 p(\lambda),
 \label{eq:hierarchical-posterior}
\end{align}
where \(\lambda\) collects hyperparameters for naturalness, correlation
lengths, truncation scales, and nuisance populations.

The likelihood is evaluated only after all structural gates pass.  A state
that violates exact normalization, color singletness, operator completeness,
or composition exclusions has zero support by construction rather than a
large chi-square penalty.

\subsection{Observable-level and matrix-element-level calibration}

The preferred calibration targets are raw observables or well-defined matrix
elements.  Let \(\cM_j(\Theta)\) be a local or bilocal matrix element and
\(\cO_i(\Theta)\) a process observable.  Then
\begin{equation}
 \log\cL
 =
 \sum_{j\in\cI_{\mathrm{ME}}}
 \log p(M_j^{\obs}\mid\cM_j)
 +
 \sum_{i\in\cI_{\proc}}
 \log p(O_i^{\obs}\mid\cO_i),
 \label{eq:matrix-and-process-likelihood}
\end{equation}
with joint covariance where the data are correlated.

Named TMD curves are calibration objects only when their extraction
likelihood, scheme, covariance, and data ancestry are sufficiently available.
Otherwise they remain external validation summaries.

\subsection{Sequential and modular calibration}

Computationally, the model may be developed in stages:
\begin{equation}
 p(\Theta_A\mid D_A)
 \longrightarrow
 p(\Theta_A,\Theta_B\mid D_A,D_B)
 \longrightarrow\cdots .
 \label{eq:sequential-calibration}
\end{equation}
Sequential updating is mathematically equivalent to joint Bayes only when the
same joint model is used and no information is discarded.  Modular
calibration, cut feedback, or staged emulator training may be useful, but the
departure from the joint posterior is declared and tested.

For example, a lattice matching analysis may supply
\(p(\theta_{\matchp}\mid D_{\latt})\) as an external posterior.  Importing it as
a prior is valid only when the data ancestry is disjoint from the downstream
likelihood or the shared data are modeled jointly.

\subsection{Tempering and robust likelihoods}

If a dataset is suspected of unmodeled contamination, the response is not an
undocumented likelihood weight.  Admissible options include:
\begin{itemize}
\item an explicit contamination mixture;
\item a heavy-tailed observation model;
\item a hierarchical between-experiment variance;
\item a diagnostic-only role;
\item a sensitivity analysis over a declared generalized-Bayes temperature.
\end{itemize}
Any likelihood tempering changes the inferential target and is stored in the
analysis manifest.

\subsection{Alternative microscopic models}

Suppose \(\mathfrak M_k\) are competing model classes.  Their predictive
distributions may be compared or combined:
\begin{equation}
 p_k(\widetilde y\mid\caldata)
 =
 \int p(\widetilde y\mid\Theta_k,\mathfrak M_k)
 p(\Theta_k\mid\caldata,\mathfrak M_k)\,\dd\Theta_k.
 \label{eq:model-predictive}
\end{equation}
Posterior model probabilities require meaningful model priors and marginal
likelihoods.  In an open model space, predictive stacking may be preferable:
\begin{equation}
 p_{\mathrm{stack}}(\widetilde y)
 =
 \sum_k w_k p_k(\widetilde y),
 \qquad
 w_k\geq0,\quad\sum_kw_k=1,
 \label{eq:stacking-mixture}
\end{equation}
with weights determined by protected out-of-sample predictive performance
\cite{YaoStacking2018}.

Model combination never merges incompatible internal states into one
microscopic member.  It is a mixture over complete model predictions.

\section{Identifiability, sensitivity, and parameter geometry}
\label{sec:identifiability}

\subsection{Structural and practical identifiability}

A parameter is structurally identifiable from observable family \(\cO\) if
\begin{equation}
 \cG(\Theta_1)=\cG(\Theta_2)
 \quad\Longrightarrow\quad
 \theta_a(\Theta_1)=\theta_a(\Theta_2)
 \label{eq:structural-identifiability}
\end{equation}
within the declared model.  Practical identifiability depends on finite data,
noise, nuisance parameters, discrepancy, and kinematic support.

An unidentifiable parameter is not rescued by reporting a narrow conditional
fit obtained after fixing correlated parameters.  The posterior geometry and
prior dependence must be exposed.

\subsection{Jacobian and Fisher information}

At a reference point,
\begin{equation}
 J_{ia}
 =
 \frac{\partial m_i}{\partial\theta_a}.
 \label{eq:sensitivity-jacobian}
\end{equation}
For a Gaussian likelihood with covariance \(\bm C\), the local expected Fisher
information is
\begin{equation}
 \bm I(\Theta)
 =
 \bm J\trans\bm C^{-1}\bm J
 +
 \bm I_{\mathrm{prior}}.
 \label{eq:fisher-information}
\end{equation}
Small singular values of
\(\bm C^{-1/2}\bm J\) identify weak parameter combinations.

The Jacobian must be computed through the same composition graph used for
predictions, using automatic differentiation, adjoint methods, or validated
finite differences.  Derivatives through hard gates or discrete model changes
are treated separately.

\subsection{Posterior correlations and profiles}

Diagnostics include:
\begin{itemize}
\item posterior correlation and partial-correlation matrices;
\item profile likelihood or profile posterior curves;
\item principal directions of the posterior covariance;
\item rank and condition number of local sensitivity matrices;
\item prior-to-posterior contraction by parameter block;
\item leave-one-dataset-family influence;
\item observable-to-parameter attribution.
\end{itemize}

A large posterior correlation is not automatically a defect; it may represent
a real physical degeneracy.  The requirement is that the degeneracy be visible
and that withheld predictions propagate it.

\subsection{Parameter ownership matrix}

Define
\begin{equation}
 A_{ia}
 =
 \begin{cases}
 1,& \theta_a\text{ is an ancestor of observable }i,\\
 0,& \text{otherwise}.
 \end{cases}
 \label{eq:ownership-matrix}
\end{equation}
The weighted sensitivity graph uses
\begin{equation}
 S_{ia}
 =
 \left|
 \frac{\theta_a}{s_i}
 \frac{\partial m_i}{\partial\theta_a}
 \right|,
 \label{eq:dimensionless-sensitivity}
\end{equation}
where \(s_i\) is a declared observable scale.

The ownership matrix tests whether a supposedly shared parameter actually
propagates to all required projections.  It also detects illicit latent
parameters that reach only one named TMD.

\begin{requirement}[Shared-physics sensitivity, V6.ID.1]
Every parameter associated with a quark--gluon vertex, OAM mixing, Wilson
kernel, nuclear current, or matching coefficient must have the complete
observable ancestry predicted by the formal operator graph.  Missing
sensitivities are implementation failures.
\end{requirement}

\subsection{Global sensitivity}

For nonlinear models, variance-based or derivative-based global sensitivity
may supplement the local Fisher analysis.  The sensitivity calculation is
performed under the prior or posterior distribution appropriate to the
question.  A parameter can be locally invisible at one point but globally
important through thresholds, nodes, or interference.

Global sensitivity is not used to delete a physically required operator.
Low sensitivity in the current data region may instead identify a target for
new measurements.

\subsection{Identifiability across observable families}

Cross-channel inference is central.  A parameter combination poorly
identified by one family may be separated by another:
\begin{equation}
 \bm I_{\mathrm{total}}
 =
 \sum_{f\in\text{families}}
 \bm J_f\trans\bm C_f^{-1}\bm J_f .
 \label{eq:family-information-sum}
\end{equation}
Examples include:
\begin{itemize}
\item spectra and currents separating Hamiltonian and current coefficients;
\item TMD and GPD marginals separating intrinsic transverse structure from
      normalization;
\item SIDIS and Drell--Yan separating link orientation;
\item Sivers and Boer--Mulders separating target-spin and active-quark-spin
      interference;
\item \(b_1\), elastic form factors, and tagged observables separating S--D
      structure, sea tensor polarization, and non-nucleonic sectors;
\item multiple \(Q\) values separating boundary and evolution parameters.
\end{itemize}

\section{Emulators, reduced models, and solver error}
\label{sec:emulators}

\subsection{Purpose and admissible role}

If the exact forward solver is too costly for posterior sampling, an emulator
\begin{equation}
 \widehat{\cG}(\Theta)
 =
 \cG(\Theta)+\delta_{\emul}(\Theta)
 \label{eq:emulator-definition}
\end{equation}
may replace selected deterministic nodes.  The emulator approximates the
declared solver; it does not replace physical discrepancy.

Admissible methods include Gaussian-process response surfaces, polynomial or
sparse-grid approximations, eigenvector continuation, reduced-basis
variational emulators, and neural surrogates.  Their errors and domains differ
and must be tested independently
\cite{Frame2018,ZhangFurnstahl2022,Drischler2021}.

\subsection{Training, validation, and test designs}

Parameter points are partitioned into:
\begin{equation}
 \cX_{\emul}
 =
 \cX_{\mathrm{train}}
 \,\dot\cup\,
 \cX_{\mathrm{val}}
 \,\dot\cup\,
 \cX_{\mathrm{test}}.
 \label{eq:emulator-split}
\end{equation}
The test set is not used for hyperparameter tuning.  It should include:
\begin{itemize}
\item space-filling points in the prior-relevant domain;
\item posterior-relevant points found adaptively;
\item edges and near-singular regions;
\item regulator and truncation transitions;
\item points near nodes, thresholds, and interference sign changes.
\end{itemize}

\subsection{Emulator covariance}

For probabilistic emulators,
\begin{equation}
 \widehat{\cG}(\Theta)
 \sim
 \mathcal N\!\left(
 \bm\mu_{\emul}(\Theta),
 \bm C_{\emul}(\Theta)
 \right).
 \label{eq:probabilistic-emulator}
\end{equation}
The emulator covariance is added to the likelihood or sampled as a latent
error only after calibration against exact solver residuals.

For deterministic reduced-basis emulators, an a posteriori error bound or
validated empirical residual plays the same role.  A small training error is
not sufficient.

\subsection{Observable and derivative fidelity}

An emulator is validated for:
\begin{itemize}
\item each observable family used in the likelihood;
\item cross-observable covariance;
\item gradients or adjoints used by HMC or sensitivity analysis;
\item exact constraints and sum rules;
\item sign changes and zero limits;
\item posterior predictive tails.
\end{itemize}
A surrogate that reproduces central values but violates a Ward identity or
gives inaccurate derivatives is not acceptable for inference.

\subsection{Adaptive refinement}

During inference, define a risk score
\begin{equation}
 R_{\emul}(\Theta)
 =
 \frac{\norm{\bm C_{\obs}^{-1/2}
 [\cG(\Theta)-\widehat{\cG}(\Theta)]}}
 {1+\norm{\bm C_{\obs}^{-1/2}\cG(\Theta)}}.
 \label{eq:emulator-risk}
\end{equation}
High posterior mass at large predicted emulator risk triggers exact solver
evaluations and retraining.  The retraining sequence is versioned, and final
validation uses an independent exact test set.

\begin{requirement}[No silent extrapolation, V6.EMUL.1]
An emulator prediction outside its certified domain is marked unavailable or
forces an exact calculation.  Extrapolation is not assigned zero emulator
uncertainty.
\end{requirement}

\subsection{Multi-fidelity hierarchy}

Let \(\cG_\ell\) denote a fidelity level.  A multi-fidelity model may use
\begin{equation}
 \cG_H(\Theta)
 =
 \rho(\Theta)\cG_L(\Theta)
 +
 \Delta_{HL}(\Theta),
 \label{eq:multifidelity}
\end{equation}
but \(\Delta_{HL}\) is an emulator relation, not physical model discrepancy.
Fidelity labels may include basis size, Fock depth, Wilson order, quadrature
resolution, or perturbative order only when the mapping is mathematically
defined.

\section{Computational inference and algorithm validation}
\label{sec:computation}

\subsection{Algorithm selection}

Gradient-based Hamiltonian Monte Carlo and \NUTS\ are natural for continuous,
differentiable, moderately high-dimensional posteriors
\cite{HoffmanGelman2014}.  Sequential Monte Carlo, nested sampling, reversible
jump methods, or tempered transitions may be required for multimodal or
discrete model spaces.  Variational or Laplace approximations are allowed only
with direct validation against exact or high-quality sampling on representative
subproblems.

The algorithm is selected after examining posterior geometry, not by a fixed
project-wide preference.

\subsection{Synthetic recovery}

For known synthetic parameters
\begin{equation}
 \Theta_0\sim p(\Theta),
 \qquad
 y_0\sim p(y\mid\Theta_0),
 \label{eq:synthetic-generation}
\end{equation}
the complete pipeline must recover the appropriate posterior uncertainty.
Tests include:
\begin{itemize}
\item one-parameter and correlated multi-parameter recovery;
\item nuisance and physical parameter separation;
\item discrepancy-on and discrepancy-off cases;
\item active/inactive Fock and Wilson sectors;
\item nonidentifiable parameter combinations;
\item emulator and exact-solver comparison.
\end{itemize}

\subsection{Simulation-based calibration}

For repeated prior-predictive simulations, \SBC\ rank statistics test whether
the inference algorithm samples the implemented posterior
\cite{Talts2018}.  For scalar test quantity \(h(\Theta)\), define
\begin{equation}
 r_s
 =
 \sum_{m=1}^{M}
 \bbone\!\left[
 h(\Theta_{s,m}^{\post})
 <
 h(\Theta_s^{\mathrm{true}})
 \right].
 \label{eq:sbc-rank}
\end{equation}
Under calibrated computation, ranks are discrete uniform after accounting for
ties and autocorrelation.

SBC validates the implementation of the probabilistic program and algorithm;
it does not validate the scientific model against nature.

\subsection{MCMC diagnostics}

For every continuous parameter and important derived quantity, report:
\begin{itemize}
\item rank-normalized split \(\Rhat\);
\item bulk and tail effective sample size;
\item Monte Carlo standard errors;
\item divergences and energy diagnostics for HMC;
\item chain and rank plots;
\item sensitivity to initialization and adaptation;
\item mode discovery or evidence of nonmixing.
\end{itemize}
Rank-normalized and localized diagnostics are used for heavy-tailed or
nonstationary chains \cite{VehtariRhat2021}.

No single numerical threshold proves convergence.  Diagnostics are combined
with repeated runs and problem-specific quantities.

\subsection{Approximate inference}

If \(q(\Theta)\) approximates the posterior, assess:
\begin{equation}
 \E_q[h(\Theta)],
 \qquad
 \Var_q[h(\Theta)],
 \qquad
 p_q(\cY_{\hold}),
 \label{eq:approx-inference-tests}
\end{equation}
against a trusted sampler on reduced problems and importance-weighted or
simulation-based diagnostics on the full problem.  Underdispersion in
variational approximations is treated as a failure unless corrected and
validated.

\subsection{Reproducibility manifest}

Every inferential run stores:
\begin{itemize}
\item code, data, formalism, and configuration hashes;
\item random seeds and chain identifiers;
\item exact and emulator solver versions;
\item prior, likelihood, discrepancy, and partition manifests;
\item sampler and adaptation settings;
\item posterior draw checksums;
\item hardware and numerical precision;
\item all warnings, rejected states, and failed diagnostics.
\end{itemize}


\section{Posterior prediction and the common microscopic ensemble}
\label{sec:posterior-prediction}

\subsection{Posterior predictive distribution}

Let \(\Theta\) denote all continuous parameters, discrete model labels,
nuisance variables, and latent discrepancy fields retained after calibration.
For a future or withheld dataset \(\cY_*\) under design and measurement record
\(\mathfrak d_*\), the posterior predictive distribution is
\begin{equation}
 p(\cY_*\mid\caldata,\mathfrak d_*)
 =
 \sum_m\int
 p(\cY_*\mid\Theta,m,\mathfrak d_*)\,
 p(\Theta,m\mid\caldata)\,
 \dd\Theta .
 \label{eq:posterior-predictive-general}
\end{equation}
The prediction includes the experimental sampling model appropriate to
\(\mathfrak d_*\).  The corresponding latent-physics predictive distribution,
with future measurement noise removed, is
\begin{equation}
 p(\cO_*\mid\caldata,\mathfrak d_*)
 =
 \sum_m\int
 \delta\!\left[
 \cO_*-\cG_*(\Theta,m;\mathfrak d_*)
 \right]\,
 p(\Theta,m\mid\caldata)\,
 \dd\Theta .
 \label{eq:latent-predictive}
\end{equation}
These distributions answer different questions and may not share one
uncertainty band label.

\begin{requirement}[Shared-member pushforward]
A posterior draw is propagated as one indivisible microscopic member through
every requested observable.  Resampling quark, gluon, proton, neutron,
Hamiltonian, Wilson, nuclear, matching, or evolution components independently
after inference is forbidden unless independence is a declared property of the
posterior factorization.
\end{requirement}

For posterior member \(s\), store
\begin{equation}
 \Theta_s
 \longmapsto
 \left\{
 |\Psi_N\rangle_s,\,
 |\Psi_D\rangle_s,\,
 W_{q,s},W_{\bar q,s},W_{g,s},\,
 U_{\gamma,s},\,
 \cZ_s,\,
 \cD_s,\,
 \cO_{1,s},\ldots,\cO_{K,s}
 \right\}.
 \label{eq:common-member-map}
\end{equation}
This member identity is the computational realization of the claim that
cross-flavor, cross-TMD, quark--gluon, proton--neutron, nucleon--deuteron, and
cross-process correlations are predictions of the common model.

\subsection{Predictive covariance and decomposition}

For latent observables \(\cO_i\), the total posterior covariance is
\begin{equation}
 \Cov_{\pred}(\cO_i,\cO_j)
 =
 \E_{\post}\!\left[
 \Cov(\cO_i,\cO_j\mid\Theta)
 \right]
 +
 \Cov_{\post}\!\left(
 \E[\cO_i\mid\Theta],
 \E[\cO_j\mid\Theta]
 \right).
 \label{eq:law-total-covariance}
\end{equation}
The first term contains stochastic discrepancy, emulator residual, and any
future-process randomness retained conditionally on \(\Theta\).  The second is
the pushforward of posterior uncertainty in shared parameters and discrete
models.

For a transparent error budget, define conditional variance contributions by
holding all but one uncertainty block fixed and using a common posterior
reference distribution:
\begin{equation}
 V_i^{(a)}
 =
 \E_{\Theta_{\neg a}}
 \Var_{\Theta_a}
 \left[
 \cO_i(\Theta_a,\Theta_{\neg a})
 \mid\Theta_{\neg a}
 \right].
 \label{eq:conditional-variance-component}
\end{equation}
Because nonlinear interactions among blocks generally exist,
\begin{equation}
 \Var(\cO_i)
 \neq
 \sum_a V_i^{(a)} .
 \label{eq:variance-not-additive}
\end{equation}
The report therefore includes interaction or remainder terms rather than
forcing all sources into a misleading quadrature sum.

\subsection{Simultaneous bands and correlation products}

Pointwise posterior intervals do not control a whole function.  For a curve
\(f(z)\) on a declared domain \(\cZ\), construct a simultaneous band from the
posterior distribution of
\begin{equation}
 T_s
 =
 \sup_{z\in\cZ}
 \frac{
 \abs{f_s(z)-\widehat f(z)}
 }{
 \widehat\sigma_f(z)+\eps_{\mathrm{floor}}
 } .
 \label{eq:simultaneous-band-statistic}
\end{equation}
If \(c_{1-\alpha}\) is the posterior quantile of \(T_s\), report
\begin{equation}
 \widehat f(z)
 \pm
 c_{1-\alpha}\widehat\sigma_f(z)
 \label{eq:simultaneous-band}
\end{equation}
together with the domain, grid or function-space representation, and floor
used in \cref{eq:simultaneous-band-statistic}.

The principal covariance deliverables include:
\begin{itemize}
\item full or compressed covariance matrices for named observable vectors;
\item correlation matrices across flavor, target polarization, parton species,
nuclear mechanism, scale, and process;
\item principal covariance modes and the parameter blocks that generate them;
\item cross-covariance needed for ratios, asymmetries, and combined fits;
\item posterior member files sufficient to reconstruct nonlinear derived
quantities without Gaussian approximation.
\end{itemize}

\subsection{Posterior predictive checks}

For a discrepancy statistic \(T(\cY,\Theta)\), generate replicated data
\(\cY^{\repdata}\) and compare
\begin{equation}
 T(\cY^{\repdata},\Theta)
 \quad\hbox{with}\quad
 T(\cY^{\obs},\Theta).
 \label{eq:ppc}
\end{equation}
Checks are chosen before inspecting their final result and include:
\begin{itemize}
\item residual shape and correlation versus \(x,Q,\qT,t,\phi\), target state,
and process;
\item normalization and shape summaries separately;
\item symmetry, support, positivity, number, momentum, current, and angular
condition residuals;
\item extreme-tail, node, sign, and zero-crossing statistics;
\item within-family and cross-family correlations;
\item replicate behavior after the same acceptance, binning, unfolding, and
derived-observable procedure used for the data.
\end{itemize}

A posterior predictive \(p\)-value is a model-criticism diagnostic, not a
frequentist significance level.  It is reported with the statistic and the
reference distribution rather than converted into a universal pass/fail
number.

\subsection{Conditional predictive checks}

A global fit may hide a failure in a weakly weighted family.  Define
conditional predictive distributions such as
\begin{equation}
 p(
 \cY_{\mathrm{tensor}}
 \mid
 \cY_{\mathrm{spectroscopy}},
 \cY_{\mathrm{current}},
 \cY_{\mathrm{unpolarized}}
 )
 \label{eq:conditional-predictive-family}
\end{equation}
or
\begin{equation}
 p(
 \cY_{\mathrm{BM}}
 \mid
 \cY_{\mathrm{Sivers}},
 \cY_{\mathrm{state}},
 \cY_{\mathrm{current}}
 ).
 \label{eq:sivers-to-bm-predictive}
\end{equation}
These tests isolate the information flow that justifies a claimed prediction.

\section{Withheld validation and falsifiability}
\label{sec:withheld-validation}

\subsection{Three distinct uses of data}

The inference program separates, as part of an iterative but auditable Bayesian workflow \cite{GelmanWorkflow2020}:
\begin{enumerate}
\item \textbf{calibration data} \(\caldata\), which update the posterior;
\item \textbf{diagnostic data} \(\diagdata\), which may guide model criticism
but do not enter the final posterior without a new versioned analysis;
\item \textbf{withheld data} \(\holddata\), whose values or final comparison
remain inaccessible until the model, priors, discrepancy, and acceptance
criteria are frozen.
\end{enumerate}
Moving a failed withheld family into calibration creates a new model version
and requires a new independent holdout.  It cannot preserve the original
prediction claim.

\subsection{Strong holdout units}

Randomly withholding individual points is often too weak because neighboring
kinematics, shared systematics, and the same extraction assumptions leak
information.  Preferred holdout units are:
\begin{itemize}
\item an entire observable family;
\item an experiment or run period;
\item a disjoint kinematic region;
\item a polarization or flavor channel;
\item a process with a different Wilson-link or hard-factor structure;
\item a lattice ensemble or operator family;
\item a nuclear observable not used in the nucleon calibration;
\item a higher Fock, Wilson, or perturbative order benchmark.
\end{itemize}

\begin{definition}[Cross-channel prediction]
A cross-channel prediction is a posterior predictive distribution for an
observable whose normalization and shape were not directly calibrated and
whose dominant operator projector, target channel, process map, or nuclear
mechanism differs from the calibration family, while sharing identifiable
microscopic parameters with it.
\end{definition}

\subsection{Predeclared validation record}

Each withheld gate has a versioned record
\begin{equation}
 \mathfrak v
 =
 \left(
 \begin{gathered}
 \text{id},\ \text{family},\ \text{data ancestry},\
 \text{information embargo},\\
 \text{prediction object},\ \text{metric},\ \text{tolerance},\
 \text{decision rule},\\
 \text{permitted revision classes}
 \end{gathered}
 \right).
 \label{eq:validation-record}
\end{equation}
The record is committed before unblinding.  Tolerances are justified by the
scientific use of the prediction, the estimated uncertainties, and analytic or
simulation calibration; they are not chosen to make an observed result pass.

\subsection{Predictive metrics}

No single metric is sufficient.  The standard validation report includes:
\begin{align}
 \mathrm{LogScore}
 &=
 -\log p(y_{\hold}\mid\caldata),
 \label{eq:log-score}\\
 \mathrm{Mahalanobis}^2
 &=
 (\bm y_{\hold}-\bm\mu_{\pred})^{\mathsf T}
 \bm\Sigma_{\pred}^{-1}
 (\bm y_{\hold}-\bm\mu_{\pred}),
 \label{eq:mahalanobis-predictive}\\
 \mathrm{Coverage}_{\alpha}
 &=
 \frac{1}{N_{\hold}}
 \sum_{i=1}^{N_{\hold}}
 \bbone[
 y_i\in C_{i,\alpha}^{\pred}
 ],
 \label{eq:predictive-coverage}\\
 \mathrm{PIT}_i
 &=
 F_i^{\pred}(y_i),
 \label{eq:pit}
\end{align}
together with family-appropriate proper scores, sign or node recovery,
integrated shape distances, and correlated residual plots.  Strictly proper
scores reward honest predictive distributions rather than only narrow central
curves \cite{GneitingRaftery2007}.  Calibration is assessed together with
sharpness \cite{GneitingCalibration2007}.

When a covariance is singular because exact constraints remove directions,
\cref{eq:mahalanobis-predictive} is evaluated in the supported subspace with
the rank and projector reported.

\subsection{Family-wise decision rule}

A withheld family passes only if:
\begin{enumerate}
\item its data ancestry is disjoint from calibration at the declared level;
\item the computational and posterior diagnostics have passed;
\item no post-unblinding parameter or discrepancy adjustment occurred;
\item all primary metrics satisfy their predeclared ranges;
\item no exact structural constraint fails;
\item residuals show no physically organized failure that is hidden by a
global scalar score;
\item the result is stable under the required prior, discrepancy, emulator,
and numerical sensitivity analyses.
\end{enumerate}
A partial pass may be reported only with the failed sub-gates and a restricted
claim.

\subsection{Canonical cross-channel tests}

The program prioritizes tests in which one physical sector controls multiple
observables:
\begin{enumerate}
\item calibrate spectroscopy and local currents, then predict flavor-resolved
PDF or form-factor information;
\item calibrate selected unpolarized and helicity information, then predict
transversity, OAM, or GTMD moments not used in the fit;
\item constrain a common quark rescattering kernel with selected Sivers data,
then predict Boer--Mulders, tensor naive-\(T\)-odd, or antiquark channels;
\item constrain the same gauge sector and state with quark data, then predict
the hierarchy or sign structure of independent gluon \(f\)- and \(d\)-type
channels without process weights;
\item calibrate nucleon properties and long-distance deuteron observables,
then predict \(b_1\), tagged tensor structure, or deuteron GTMD moments;
\item calibrate low-\(q_T\) data in one established factorization channel,
then predict another process after applying the independently specified link,
hard, and color map;
\item calibrate lower truncation levels, then predict the next basis, Fock,
Wilson, or perturbative order.
\end{enumerate}

\subsection{Failure taxonomy}

A failed withheld family is classified by tracing the result through the
typed graph:
\begin{longtable}{L{0.22\textwidth}L{0.32\textwidth}L{0.37\textwidth}}
\toprule
Class & Diagnostic signature & Permitted response\\
\midrule
\endhead
State deficiency &
Multiple reductions fail coherently; missing flavor, helicity, OAM, or Fock
support. &
Revise the Hamiltonian, retained sectors, regulator, or state truncation; rerun
state and current calibration.\\
Current/operator deficiency &
Spectroscopy closes but local moments or selected projections fail. &
Add or rematch a current/operator consistently with the Hamiltonian and all
affected reductions.\\
Wilson deficiency &
Link-even quantities close but link-odd sign, rank, or process reversal fails. &
Revise intermediate states, pole prescriptions, Wilson order, soft separation,
or color contractions.\\
Nuclear deficiency &
Nucleon tests pass while binding, \(b_1\), tagged, tensor, or coherent families
fail. &
Revise the matched nuclear Hamiltonian, currents, sectors, or coherent
amplitudes; preserve nucleon member identity.\\
Matching/evolution deficiency &
Low-resolution moments close but scale or small-\(b\) behavior fails. &
Revise the closed operator basis, LF-to-QCD matching, anomalous dimensions,
threshold maps, or accuracy manifest.\\
Process deficiency &
Universal objects close but a process prediction fails. &
Reassess hard factors, fragmentation, link/color map, power corrections,
Glauber status, or factorization claim.\\
Discrepancy deficiency &
Residual structure is compatible with a known omitted order but not with the
declared kernel. &
Revise the physically owned discrepancy basis and repeat calibration with a
new version.\\
Data/likelihood deficiency &
Failure tracks data ancestry, covariance, acceptance, or extraction choices. &
Correct the likelihood or data object; do not change the microscopic model
before the data issue is resolved.\\
\bottomrule
\end{longtable}

\begin{requirement}[No curve-specific repair]
A failure may not be repaired by adding a coefficient, width, normalization,
or latent curve attached only to the failed named TMD.  A new parameter must
belong to a physical interaction, state sector, current, Wilson kernel,
matching coefficient, process term, or explicit discrepancy operator and must
propagate to every object that contains that source.
\end{requirement}

\subsection{Unblinding and revision protocol}

The unblinding record contains:
\begin{enumerate}
\item the frozen model and data hashes;
\item the original posterior draw hashes;
\item the withheld predictions generated before access;
\item the time and identity of unblinding;
\item the automatic gate output;
\item the scientific interpretation written before any refit;
\item the revision decision and the claim retired or retained.
\end{enumerate}
After revision, the original prediction remains archived.  A new analysis
receives a new semantic version and a new holdout.

\section{Predictive comparison, robustness, and influence}
\label{sec:model-comparison}

\subsection{Cross-validation unit}

Pointwise leave-one-out cross-validation is useful only when one observation
is a legitimate conditionally independent predictive unit.  Shared
normalizations, replicas, kinematic curves, or common events require
leave-group-out units.  Define groups \(g=1,\ldots,G\) by experiment,
systematic block, curve, observable family, or data ancestry:
\begin{equation}
 \mathrm{elpd}_{\LOO}
 =
 \sum_{g=1}^{G}
 \log
 p(
 \bm y_g
 \mid
 \bm y_{-g}
 ).
 \label{eq:group-loo}
\end{equation}
Pareto-smoothed importance sampling may approximate refits and provides a
tail-shape diagnostic, but influential groups with unreliable importance
weights are refit directly \cite{VehtariLOO2017,VehtariPSIS2024}.

Cross-validation estimates predictive performance for the represented data
distribution.  It does not replace a scientifically distinct withheld process
or observable family.

\subsection{Predictive stacking}

For candidate model variants \(M_k\), predictive stacking selects
\(w_k\geq0\), \(\sum_k w_k=1\), to maximize held-out predictive score:
\begin{equation}
 \widehat{\bm w}
 =
 \arg\max_{\bm w}
 \sum_g
 \log\left[
 \sum_k
 w_k\,
 p_k(\bm y_g\mid\bm y_{-g})
 \right].
 \label{eq:stacking-optimization}
\end{equation}
Stacking is a predictive combination of declared alternatives, not permission
to add mutually exclusive mechanisms in one microscopic state
\cite{YaoStacking2018}.  Each stacked component retains its own physical
identity and provenance.

Bayes factors or posterior model probabilities are reported only when the
candidate set and prior mass assignment make the closed-model assumption
scientifically defensible.  Otherwise, predictive comparison and stacking are
preferred.

\subsection{Prior sensitivity}

For each load-bearing parameter block, report:
\begin{itemize}
\item prior-to-posterior contraction;
\item sensitivity to alternative physically credible prior scales and shapes;
\item changes in withheld predictive distributions;
\item prior predictive pathologies;
\item whether an apparently constrained parameter is informed by data or by an
exact/truncation gate.
\end{itemize}
A parameter is not described as ``measured'' when its posterior is mainly a
restatement of the prior.

\subsection{Discrepancy sensitivity}

Repeat the analysis across a finite set of physically justified discrepancy
bases, correlation lengths, and power-counting hyperpriors.  Monitor:
\begin{equation}
 \Delta\theta_a,
 \qquad
 \Delta p(\cY_{\hold}),
 \qquad
 \Delta\mathrm{elpd},
 \qquad
 \Delta \Cov(\cO_i,\cO_j).
 \label{eq:discrepancy-sensitivity}
\end{equation}
Strong variation signals confounding between physical parameters and model
discrepancy.  Flexible discrepancy priors can improve interpolation while
destroying physical parameter identifiability; physically meaningful priors
and orthogonality constraints are therefore indispensable
\cite{KennedyOHagan2001,BrynjarsdottirOHagan2014}.

\subsection{Influence analysis}

For data group \(g\), define an influence measure using the change in a
scientific target \(h(\Theta)\):
\begin{equation}
 I_g[h]
 =
 \frac{
 \abs{
 \E[h(\Theta)\mid\allData]
 -
 \E[h(\Theta)\mid\allData\setminus\bm y_g]
 }
 }{
 \sqrt{\Var[h(\Theta)\mid\allData]}+\eps
 } .
 \label{eq:group-influence}
\end{equation}
Report high-influence experiments, lattice ensembles, extracted inputs, and
calibration observables together with their data ancestry and covariance
assumptions.

\subsection{Robustness to numerical and formal choices}

The final scientific conclusion must be stable, within its reported
uncertainty, under the declared variations in:
\begin{itemize}
\item quadrature, grid, and transform resolution;
\item solver tolerances and eigensolver state tracking;
\item emulator architecture and training design;
\item basis, Fock, Wilson, and perturbative truncation;
\item regulator and renormalization choices with refitting;
\item factorization profile and scale paths;
\item data binning, covariance regularization, and nuisance representation;
\item posterior sampler and parameterization.
\end{itemize}
These are not all averaged together.  Each variation is tagged as an
equivalent numerical representation, a truncation axis, a scheme choice, or a
different physical model.

\section{Experiment design and information geometry}
\label{sec:design}

\subsection{Expected information gain}

Let \(d\) denote a candidate experimental design: beam energy, target state,
kinematic bin, observable, luminosity allocation, detector resolution, or
analysis strategy.  For scientific parameters or mechanism label
\(\Theta_{\mathrm{tar}}\), the expected information gain is
\begin{equation}
 \mathrm{EIG}(d)
 =
 \E_{\bm y\sim p(\bm y\mid d,\caldata)}
 \left[
 \KL\!\left(
 p(\Theta_{\mathrm{tar}}\mid\bm y,d,\caldata)
 \,\Vert\,
 p(\Theta_{\mathrm{tar}}\mid\caldata)
 \right)
 \right].
 \label{eq:eig}
\end{equation}
The target may be a physical parameter block, a prediction, a sign or node, a
factorization alternative, or a mechanism discriminator.  This is the standard
Bayesian experimental-design objective, extended here by the typed physics and
feasibility constraints \cite{ChalonerVerdinelli1995,CoonsHuan2025}.

A design that only reconstrains an already narrow direction can have low EIG
even when its projected statistical error is small.

\subsection{Fisher approximation}

Near a regular reference point with Gaussian noise covariance
\(\bm\Sigma_y(d)\), a local approximation is
\begin{equation}
 \bm F_d
 =
 \bm J_d^{\mathsf T}
 \bm\Sigma_y(d)^{-1}
 \bm J_d
 +
 \bm F_{\mathrm{prior}},
 \qquad
 (\bm J_d)_{ia}
 =
 \frac{\partial \cO_i(d)}{\partial\theta_a}.
 \label{eq:fisher-design}
\end{equation}
The Fisher approximation is accepted only after comparison with posterior or
simulation-based design calculations in regions with nonlinearity,
multimodality, boundaries, or discrete mechanisms.

\subsection{Mechanism discrimination}

For competing mechanisms \(M_1,M_2\), define a predictive separation:
\begin{equation}
 \cD_{\mathrm{sep}}(d)
 =
 \E_{\bm y\sim p_1}
 \log\frac{p_1(\bm y\mid d)}{p_2(\bm y\mid d)}
 +
 \E_{\bm y\sim p_2}
 \log\frac{p_2(\bm y\mid d)}{p_1(\bm y\mid d)} .
 \label{eq:mechanism-separation}
\end{equation}
This symmetric predictive divergence identifies observables that distinguish
physical mechanisms rather than only improve one parameter variance.

Candidate applications include:
\begin{itemize}
\item Sivers versus Boer--Mulders spin projectors under a common rescattering
kernel;
\item gluon \(f\)- versus \(d\)-type color channels;
\item S--D versus non-nucleonic tensor mechanisms;
\item impulse versus coherent small-\(x\) contributions;
\item straight- versus staple-link OAM observables;
\item competing LF-to-QCD matching or nuclear EFT truncations.
\end{itemize}

\subsection{Feasibility-aware utility}

The design objective includes cost, acceptance, systematic risk, and
factorization status:
\begin{equation}
 U(d)
 =
 \mathrm{EIG}(d)
 -
 \lambda_C C(d)
 -
 \lambda_S R_{\sys}(d)
 -
 \lambda_F R_{\mathrm{fact}}(d).
 \label{eq:design-utility}
\end{equation}
A high-information observable in a process with unresolved color entanglement
does not automatically outrank a slightly less informative observable in an
established factorization channel.

\subsection{Sequential design}

After a pilot dataset \(\bm y_1\), update the posterior and choose
\begin{equation}
 d_2^*
 =
 \arg\max_{d_2}
 \E[
 U(d_2)
 \mid
 \bm y_1,d_1,\caldata
 ].
 \label{eq:sequential-design}
\end{equation}
The design history, interim stopping rules, and any adaptive choice are
recorded so the final likelihood and validation partition remain correct.

When the high-fidelity forward model is expensive, multi-fidelity estimators
or validated emulators may accelerate EIG calculations, but the final design
ranking must include emulator uncertainty and spot checks against the exact
solver.

\section{Reporting and uncertainty communication}
\label{sec:reporting}

\subsection{Minimum result package}

Every scientific release includes:
\begin{enumerate}
\item the posterior member ensemble or a loss-controlled compressed
representation;
\item central summaries and pointwise plus simultaneous intervals;
\item covariance and correlation products;
\item separated uncertainty blocks and nonlinear interaction remainder;
\item calibration, diagnostic, and withheld data manifests;
\item prior, likelihood, discrepancy, emulator, and sampler manifests;
\item posterior predictive and withheld validation reports;
\item exact hashes of code, data, formalism, and generated artifacts;
\item evidence labels and factorization/matching status for every result.
\end{enumerate}

\subsection{Evidence vocabulary}

A released result carries one of the following evidence statuses:
\begin{longtable}{L{0.18\textwidth}L{0.72\textwidth}}
\toprule
Status & Meaning\\
\midrule
\endhead
Exact &
Consequence of operator definition, representation, symmetry, kinematics,
algebraic composition, or a proven theorem under stated hypotheses.\\
Calibrated &
Directly updated by the declared calibration likelihood; posterior uncertainty
and prior sensitivity are reported.\\
Lattice-informed &
Updated by matched lattice matrix elements or an explicit lattice prior with
scheme and covariance recorded.\\
Predicted--withheld &
Normalization and shape excluded from calibration and successfully tested
against a frozen withheld family.\\
Predicted--unobserved &
Posterior prediction not directly calibrated, but no relevant independent
measurement has yet tested it.\\
Model sensitivity &
Physically allowed but underdetermined alternative, used to expose model
dependence rather than a calibrated probability statement.\\
Temporary &
Controlled implementation boundary with a named replacement and distinguishing
test.\\
Exploratory &
Process, matching, factorization, or data status is not sufficient for an
authoritative physical claim.\\
\bottomrule
\end{longtable}

Passing a structural test does not upgrade a model sensitivity to a calibrated
or withheld prediction.

\subsection{Credible intervals and coverage language}

Posterior credible intervals express probability under the declared model,
prior, likelihood, and discrepancy.  They are not called confidence intervals.
Empirical coverage may be assessed by SBC, synthetic benchmarks, or repeated
withheld families, but the achieved calibration is reported separately.

Likewise, an envelope over wave functions, regulator choices, or model
alternatives is not a posterior credible interval unless those alternatives
carry a declared probability model and have been sampled accordingly.

\subsection{Central summaries}

For skewed or multimodal posteriors, report medians, highest-density regions,
or modes only with the full distribution or member ensemble.  A mean curve
assembled point by point may not correspond to any physical state.  At least
one representative posterior member near the central functional region is
identified for reproducible benchmark calculations.

\subsection{Ratios and asymmetries}

Derived ratios and asymmetries are calculated member by member:
\begin{equation}
 A_s
 =
 \frac{N_s}{D_s},
 \qquad
 R_s
 =
 \frac{\cO_{1,s}}{\cO_{2,s}},
 \label{eq:memberwise-ratio}
\end{equation}
not by dividing separately summarized numerator and denominator bands.  The
report identifies denominator-zero or sign-change regions and any non-Gaussian
tails.

\subsection{No hidden uncertainty inflation}

A mismatch may not be hidden by:
\begin{itemize}
\item adding experimental and theory errors in an undocumented way;
\item rescaling all uncertainties until reduced \(\chi^2\) is near one;
\item clipping negative eigenvalues of a covariance or correlator without
diagnosis;
\item replacing a failed predictive mean by a broad unstructured discrepancy;
\item reporting only a total band that erases which physical layer failed.
\end{itemize}
All covariance regularization, nuisance constraints, and discrepancy
hyperpriors are part of the released manifest.


\section{Software realization and formal interfaces}
\label{sec:software}

\subsection{Core objects}

The inference layer is represented by typed objects rather than one monolithic
fit script:
\begin{longtable}{L{0.25\textwidth}L{0.66\textwidth}}
\toprule
Object & Responsibility\\
\midrule
\endhead
\texttt{ParameterId} &
Stable physical ownership, units, support, symmetry, regulator, source object,
differentiability, and transformation.\\
\texttt{ParameterBlock} &
Hierarchical prior, exact constraints, shared-member identity, and allowed
couplings for one owned parameter family.\\
\texttt{DataSetId} &
Observable, experiment, event ancestry, kinematics, units, convention, scheme,
covariance, likelihood, and partition status.\\
\texttt{DataAncestryGraph} &
Source-event and extraction lineage, overlap and exclusion relations, and
double-counting rejection.\\
\texttt{LikelihoodTerm} &
Typed mapping from a prediction and nuisance structure to a normalized
probability density or mass.\\
\texttt{DiscrepancyField} &
Physically owned omitted-order or omitted-sector field with basis, kernel,
power counting, correlation domain, and orthogonality constraints.\\
\texttt{PriorModel} &
Prior and hyperprior graph, naturalness scales, symmetry-breaking hierarchy,
and prior-predictive simulator.\\
\texttt{ForwardModel} &
Versioned deterministic map from a complete microscopic member to all required
matrix elements and observables.\\
\texttt{Emulator} &
Surrogate prediction, derivative, covariance, domain, exact-training hashes,
and validation report.\\
\texttt{PosteriorProgram} &
Joint log density, gradients where available, discrete alternatives, and
sampling interface.\\
\texttt{CalibrationSplit} &
Frozen calibration, diagnostic, withheld, and embargo records with ancestry
checks.\\
\texttt{PosteriorMember} &
One inseparable draw containing physical, nuisance, discrepancy, truncation,
and model-label identity.\\
\texttt{PredictiveEnsemble} &
Memberwise pushforward to latent observables and replicated measurements,
retaining covariance and evidence status.\\
\texttt{ValidationGate} &
Withheld family, metric, threshold, unblinding record, and structured decision.\\
\texttt{ModelRevisionRecord} &
Failed gate, traced source layer, permitted revision, retired claim, semantic
version, and replacement holdout.\\
\texttt{ExperimentalDesign} &
Candidate settings, likelihood, feasibility, systematics, factorization status,
and information utility.\\
\texttt{InferenceManifest} &
Hashes, environment, seeds, sampler, diagnostics, posterior checks, and output
artifact checksums.\\
\bottomrule
\end{longtable}

\subsection{Typed inferential graph}

At a finite implementation level, define \(\Infer\) as a typed probabilistic
graph layered over the deterministic \(\Amp\), \(\Dens\), \(\Match\),
\(\Red\), and \(\Proc\) maps.  Nodes are:
\begin{equation}
 \left\{
 \Theta,\,
 \mathfrak m,\,
 \bm y,\,
 \bm\nu,\,
 \bm\delta,\,
 \cO,\,
 \cY^{\repdata}
 \right\},
 \label{eq:infer-nodes}
\end{equation}
and edges are conditional distributions or deterministic maps.  A composition
is valid only when species, target, operator, process, scheme, member,
kinematic, and data-ancestry identities agree or an explicit adapter is
registered.

\begin{requirement}[No implicit likelihood adapter]
Changing units, tensor convention, bin averaging, acceptance, radiative
correction, scheme, or covariance representation requires a versioned adapter.
A likelihood may not consume a numerically compatible array whose physical
identity is different.
\end{requirement}

\subsection{Stable diagnostic taxonomy}

Every fail-closed error carries:
\begin{equation}
 \mathfrak e
 =
 \left(
 \text{requirement id},
 \text{object ids},
 \text{expected type},
 \text{received type},
 \text{source path},
 \text{suggested physical repair}
 \right).
 \label{eq:diagnostic-schema}
\end{equation}
Diagnostics distinguish:
\begin{itemize}
\item exact structural mismatch;
\item missing physical input;
\item computational nonconvergence;
\item emulator-domain violation;
\item posterior nonidentifiability;
\item withheld-validation failure;
\item data-ancestry overlap;
\item matching or factorization incompleteness.
\end{itemize}
A failed structural requirement is never converted into a large variance.

\subsection{Automatic differentiation and adjoints}

Where the forward chain is differentiable, expose:
\begin{equation}
 \frac{\partial\cO_i}{\partial\theta_a},
 \qquad
 \frac{\partial\log p}{\partial\theta_a},
 \qquad
 \bm v^{\mathsf T}\bm J,
 \qquad
 \bm J\bm u .
 \label{eq:ad-contract}
\end{equation}
Gradients are validated by complex-step or high-order finite differences on
analytic benchmarks and by exact identities for normalization, symmetry, and
sum rules.  A gradient whose relative error exceeds its declared tolerance is
not used for HMC, design, or Fisher analysis.

\subsection{Persistent storage}

The ensemble store is append-only and content addressed.  It contains:
\begin{itemize}
\item raw unconstrained and transformed parameters;
\item discrete model labels;
\item per-member state, operator, matching, evolution, and process hashes;
\item log prior, log likelihood by dataset family, and log posterior;
\item sampler diagnostics and weights;
\item predicted matrix elements and observables or reproducible lazy recipes;
\item discrepancy and emulator draws;
\item evidence, matching, evolution, and factorization status.
\end{itemize}
Compression is accepted only after demonstrating preservation of target
moments, covariance, tails, multimodality, and withheld predictive scores.

\section{Validation ladder and property-based tests}
\label{sec:validation-ladder}

\subsection{Layered validation}

The inference architecture is validated in increasing depth:
\begin{enumerate}
\item \textbf{Type layer:} physical units, conventions, operator identity,
data ancestry, partition disjointness, and parameter ownership.
\item \textbf{Probability layer:} normalized likelihoods, proper priors,
positive-semidefinite covariance, Jacobians, and support.
\item \textbf{Implementation layer:} synthetic recovery, \SBC, gradients,
sampler diagnostics, and deterministic reproducibility.
\item \textbf{Forward layer:} exact-versus-emulated agreement and all
Volumes~I--V structural gates.
\item \textbf{Calibration layer:} prior predictive adequacy, residual checks,
identifiability, and sensitivity to nuisance/discrepancy choices.
\item \textbf{Predictive layer:} family-wise cross-validation, posterior
predictive calibration, and frozen withheld gates.
\item \textbf{Scientific layer:} successful cross-channel predictions and
traceable model revision after failures.
\end{enumerate}

\subsection{Mandatory property tests}

At minimum, the software must test:
\begin{enumerate}
\item every parameter has exactly one physical owner or is an explicit nuisance
or discrepancy quantity;
\item no named-TMD-only free parameter exists;
\item exact constraints map the parameter support into the physical state
space;
\item calibration, diagnostic, and withheld records are pairwise disjoint under
the ancestry graph;
\item raw data and derived extractions from the same events cannot both enter
without a joint likelihood;
\item covariance matrices are symmetric positive semidefinite in their declared
support and preserve exact null directions;
\item nuisance marginalization agrees with analytic results in linear-Gaussian
benchmarks;
\item likelihood normalization is verified numerically or analytically;
\item prior predictive samples obey structural physics gates;
\item synthetic posterior recovery returns the generating parameters within
calibrated rank behavior;
\item \SBC\ ranks are compatible with uniformity for analytic test models;
\item automatic derivatives agree with independent numerical derivatives;
\item exact and emulator predictions agree within the emulator contract;
\item emulator covariance covers held-out exact residuals at its declared
frequency or posterior level;
\item multiple chains pass rank-normalized \(\Rhat\), bulk/tail \(\ESS\), and
Monte Carlo error criteria;
\item no HMC divergences or hidden mode exclusions remain in accepted runs;
\item posterior member identity is preserved across all reductions;
\item memberwise ratios differ from ratios of separately summarized bands in a
known non-Gaussian test;
\item the law of total covariance is reproduced by Monte Carlo;
\item pointwise and simultaneous intervals are distinguished;
\item grouped \(\LOO\) matches direct refits on benchmark groups;
\item \PSIS\ tail diagnostics trigger direct refits when required;
\item stacking weights reproduce a known predictive mixture benchmark;
\item prior, discrepancy, and numerical sensitivity runs are deterministic and
traceable;
\item withheld gates cannot be evaluated before the unblinding token is
released;
\item a post-unblinding model change invalidates the old prediction claim;
\item a curve-specific repair is rejected;
\item experimental-design utilities include cost, systematics, and
factorization gates;
\item all reports can be regenerated from hashes and manifests.
\end{enumerate}

\subsection{Mandatory negative injections}

The test suite deliberately injects:
\begin{enumerate}
\item duplicate physical ownership for one parameter;
\item one parameter attached independently to two named TMDs;
\item an unowned free normalization;
\item an exact symmetry represented as a soft prior;
\item a covariance with a negative supported eigenvalue;
\item silent removal of a covariance null direction;
\item double counting of raw and extracted data;
\item overlap between calibration and withheld event ancestry;
\item use of the withheld value in emulator training;
\item an incorrect likelihood Jacobian;
\item a nonnormalized likelihood;
\item an unrecorded unit or tensor-convention conversion;
\item a discrepancy field with support identical to a physical parameter
direction and no identifiability control;
\item independent resampling of quark and gluon parts of one posterior member;
\item an emulator queried outside its certified domain;
\item a biased emulator whose covariance is too small;
\item a wrong gradient;
\item an MCMC chain trapped in one of two known modes;
\item a variational posterior with hidden tail undercoverage;
\item pointwise bands mislabeled as simultaneous;
\item a ratio formed from separate medians;
\item pointwise \(\LOO\) applied to correlated curve blocks;
\item unreliable \PSIS\ weights accepted without refit;
\item stacking of alternatives that violate a provenance exclusion;
\item a failed withheld family moved silently into calibration;
\item a new TMD-specific coefficient added after unblinding;
\item a theory-error band inflated until the gate passes;
\item an experimental design selecting a factorization-broken process without
penalty;
\item a report missing code, data, posterior, or formalism hashes.
\end{enumerate}

\section{Analytic and synthetic benchmark suite}
\label{sec:benchmarks}

\subsection{Benchmark A: linear-Gaussian shared parameter}

Let
\begin{equation}
 \bm y
 =
 \bm A\bm\theta+\bm\eps,
 \qquad
 \bm\eps\sim\cN(\bm0,\bm\Sigma_y),
 \qquad
 \bm\theta\sim\cN(\bm\mu_0,\bm\Sigma_0).
 \label{eq:benchmark-a}
\end{equation}
The exact posterior is
\begin{align}
 \bm\Sigma_{\post}^{-1}
 &=
 \bm\Sigma_0^{-1}
 +
 \bm A^{\mathsf T}
 \bm\Sigma_y^{-1}
 \bm A,
 \label{eq:benchmark-a-cov}\\
 \bm\mu_{\post}
 &=
 \bm\Sigma_{\post}
 \left(
 \bm\Sigma_0^{-1}\bm\mu_0
 +
 \bm A^{\mathsf T}\bm\Sigma_y^{-1}\bm y
 \right).
 \label{eq:benchmark-a-mean}
\end{align}
One shared \(\theta_a\) controls two observable families.  The benchmark checks
posterior, predictive covariance, nuisance marginalization, gradients,
\(\HMC\), and member identity.

\subsection{Benchmark B: parameter--discrepancy confounding}

Use
\begin{equation}
 y(x)
 =
 \theta_0+\theta_1x+\delta(x)+\eps(x),
 \label{eq:benchmark-b}
\end{equation}
with \(\delta\) drawn from:
\begin{enumerate}
\item an unrestricted smooth \(\GP\);
\item an operator-orthogonal basis;
\item a power-counting-informed correlated basis.
\end{enumerate}
The unrestricted case must demonstrate weak or biased physical parameter
identification.  The constrained cases must recover the known synthetic
parameters within their declared prior and coverage conditions.  This
benchmark tests that better interpolation is not confused with better
microscopic inference.

\subsection{Benchmark C: correlated EFT truncation}

Generate order-by-order curves
\begin{equation}
 \cO_k(z)
 =
 \cO_{\mathrm{ref}}(z)
 \sum_{n=0}^{k}
 c_n(z)Q(z)^n,
 \qquad
 c_n\sim\GP(0,K_c).
 \label{eq:benchmark-c}
\end{equation}
Withhold order \(k+1\).  The inference must recover the coefficient scale and
correlation length sufficiently to give calibrated predictive coverage of the
next order.  Independent pointwise errors are included as a failing control.

\subsection{Benchmark D: common-parent multi-observable closure}

Construct an analytic state with shared parameters
\[
 \Theta_{\mathrm{toy}}
 =
 (\theta_{\mathrm{spin}},
  \theta_{\mathrm{OAM}},
  \theta_{\mathrm{sea}},
  \theta_g,
  \theta_{\mathrm{nuc}}).
\]
Generate synthetic TMD, GPD, PDF, current, and tensor observables from one
common overlap.  Calibrate a restricted subset and withhold:
\begin{itemize}
\item a different reduction;
\item a different flavor;
\item the gluon momentum moment;
\item a tensor deuteron observable.
\end{itemize}
Independent curve normalizations can fit calibration but must fail at least
one holdout; the shared model must predict all generated holdouts within the
predeclared gate.

\subsection{Benchmark E: emulator honesty}

Train an emulator on a sparse exact design with one deliberately underresolved
region.  The benchmark requires:
\begin{itemize}
\item the mean residual to be covered by emulator covariance;
\item an out-of-domain query to fail closed;
\item posterior results with the emulator to agree with exact-solver inference
within Monte Carlo and emulator uncertainty;
\item adaptive refinement to select the underresolved region;
\item derivatives to pass an independent check.
\end{itemize}
Eigenvector-continuation and variational emulators are useful exemplars because
they preserve a physical reduced subspace rather than fitting only output
curves \cite{Frame2018,ZhangFurnstahl2022}.

\subsection{Benchmark F: shared rescattering prediction}

Define one link-odd kernel parameter block \(\theta_{\wil}\) and two distinct
projectors:
\begin{equation}
 F_{\mathrm{Siv}}
 =
 \cP_{\mathrm{Siv}}
 W_{\mathrm{odd}}(\theta_{\wil},\theta_{\mathrm{state}}),
 \qquad
 F_{\mathrm{BM}}
 =
 \cP_{\mathrm{BM}}
 W_{\mathrm{odd}}(\theta_{\wil},\theta_{\mathrm{state}}).
 \label{eq:benchmark-f}
\end{equation}
Calibrate selected Sivers-like synthetic data and withhold the Boer--Mulders-like
family.  The common-kernel model must predict the holdout; a model with an
independent Boer--Mulders normalization is marked as calibrated, not
predicted.  Additional members test tensor and antiquark projectors.

\subsection{Benchmark G: data ancestry and double counting}

Create raw pseudo-events, a binned cross section, and an extracted asymmetry
from the same events.  The ancestry graph must:
\begin{itemize}
\item permit the raw-event likelihood alone;
\item permit a joint derived-observable likelihood with cross-covariance;
\item reject multiplication of statistically overlapping summaries as
independent terms;
\item distinguish a genuinely independent replication.
\end{itemize}

\subsection{Benchmark H: predictive alternatives and stacking}

Generate data from a distribution outside two imperfect candidate models.
Verify that:
\begin{itemize}
\item posterior model probability is prior-sensitive under the closed-model
assumption;
\item grouped predictive stacking improves held-out proper score;
\item stacking weights do not merge mutually exclusive internal mechanisms;
\item the stacked predictive covariance includes within- and between-model
components.
\end{itemize}

\subsection{Benchmark I: experimental design}

Construct two candidate measurements:
\begin{enumerate}
\item one with small statistical error along an already constrained parameter
direction;
\item one with larger raw error but strong sensitivity to a poorly identified
mechanism.
\end{enumerate}
The EIG criterion should prefer the second when its feasibility and systematic
penalties are comparable.  A local Fisher ranking must agree only in the
linear-Gaussian limit.

\section{Staged implementation program}
\label{sec:staging}

The formal volume does not prescribe repository package numbers before the
preceding C5 APIs are known.  The implementation stages are therefore labeled
\(\mathrm{I0}\)--\(\mathrm{I6}\).

\subsection{I0: inferential type and ancestry spine}

Implement:
\begin{itemize}
\item parameter ownership and support;
\item dataset and covariance identity;
\item ancestry graph and partition records;
\item likelihood and discrepancy interfaces;
\item exact structural gates.
\end{itemize}
Use analytic likelihoods only.  No expensive microscopic inference is required.

\subsection{I1: analytic posterior and computational validation}

Implement Benchmarks A--C, including:
\begin{itemize}
\item analytic linear-Gaussian posteriors;
\item nuisance marginalization;
\item parameter--discrepancy confounding;
\item correlated truncation fields;
\item automatic differentiation;
\item \HMC/\NUTS\ and \SBC\ validation.
\end{itemize}

\subsection{I2: common-parent synthetic inference}

Connect the validation-only common overlap, sea/gluon, and Wilson pilots after
their APIs pass.  Implement Benchmark D and F without connecting to accepted
production.  Demonstrate shared-member propagation and frozen holdouts.

\subsection{I3: emulators and scalable posterior}

Implement validated state, matrix-element, and observable emulators; adaptive
training; multi-fidelity control; exact spot checks; and posterior propagation
of emulator uncertainty.

\subsection{I4: heterogeneous calibration likelihood}

Add spectroscopy, currents, form factors, selected PDFs or matrix elements,
lattice data, deuteron static properties, and restricted process observables.
Freeze data ancestry, calibration/diagnostic/holdout partitions, and primary
validation gates before final sampling.

\subsection{I5: full common posterior and withheld validation}

Run the shared posterior over all activated microscopic, Wilson, nuclear,
matching, evolution, nuisance, and discrepancy blocks.  Produce:
\begin{itemize}
\item common posterior members;
\item covariance products;
\item cross-channel posterior predictions;
\item frozen withheld validation;
\item structured revision decisions.
\end{itemize}

\subsection{I6: experimental design and living validation}

Use the accepted posterior to rank new measurements, update only through
versioned sequential-design records, and maintain an expanding set of
independent holdouts.  Old predictions remain immutable scientific artifacts.

\section{Risks, failure modes, and explicit nonclaims}
\label{sec:risks}

\begin{longtable}{L{0.22\textwidth}L{0.32\textwidth}L{0.37\textwidth}}
\toprule
Risk & Failure mode & Required safeguard\\
\midrule
\endhead
Statistical patchwork &
One normalization, width, or latent curve per TMD recreates the original
phenomenological assembly. &
Enforce physical parameter ownership and cross-observable propagation.\\
Data double counting &
Raw measurements, corrected tables, and extracted PDFs/TMDs from the same
events are multiplied as independent likelihoods. &
Use the ancestry graph, joint covariance, or choose one information level.\\
Discrepancy absorption &
A flexible discrepancy explains every residual and leaves physical parameters
unidentified. &
Use power counting, operator support, orthogonality, prior sensitivity, and
withheld channels.\\
Underestimated theory error &
Truncation and regulator uncertainty are omitted or treated as independent
pointwise noise. &
Use order-aware correlated discrepancy and next-order validation.\\
Overestimated theory error &
Unstructured theory covariance is inflated until all data pass. &
Freeze hyperpriors and validation gates; report layer-specific residuals.\\
Emulator bias &
A fast surrogate shifts the posterior while its uncertainty remains too small. &
Use exact holdouts, derivative tests, covariance calibration, and adaptive
refinement.\\
MCMC false convergence &
Chains share one mode, heavy tails, or divergent trajectories while simple
summaries appear stable. &
Use rank diagnostics, tail ESS, divergences, repeated parameterizations, and
synthetic recovery.\\
Weak holdout &
Random withheld points are predictable from neighboring points or shared
systematics. &
Hold out families, experiments, regions, processes, or operator channels.\\
Post-unblinding tuning &
A failed prediction is repaired and still described as withheld. &
Archive the original result, version the revision, and choose a new holdout.\\
Model comparison overreach &
Bayes factors are interpreted as truth probabilities in an open model class. &
Use predictive checks, grouped cross-validation, and stacking with explicit
candidate limitations.\\
Uncertainty-band misuse &
Pointwise intervals, wave envelopes, and credible bands receive the same label. &
Publish interval semantics, simultaneous bands, covariance, and member files.\\
Design myopia &
An experiment is optimized for parameter variance while ignoring systematics,
acceptance, cost, or factorization failure. &
Use feasibility-aware utility and process-status gates.\\
\bottomrule
\end{longtable}

The volume does not claim:
\begin{itemize}
\item that Bayesian probability makes an incomplete physical model correct;
\item that a Gaussian likelihood or discrepancy is universally adequate;
\item that posterior concentration proves parameter identifiability;
\item that posterior predictive agreement proves the model unique;
\item that pointwise coverage guarantees functional coverage;
\item that cross-validation replaces a protected scientific holdout;
\item that all uncertainty sources possess meaningful prior probabilities;
\item that model alternatives should always be averaged;
\item that a successful analytic or validation-only pilot is a physical
extraction;
\item that inference can begin before Volumes~I--V structural gates are
satisfied at the activated scope.
\end{itemize}

\section{Acceptance criteria and handoff}
\label{sec:acceptance}

\subsection{Formal acceptance criteria}

Volume VI is implemented faithfully only when:
\begin{enumerate}
\item every inferred parameter has a unique physical, nuisance, or discrepancy
owner;
\item exact constraints and provenance exclusions are enforced structurally;
\item dataset ancestry prevents information double counting;
\item the calibration, diagnostic, and withheld partition is frozen and
machine readable;
\item all likelihoods, covariance adapters, and nuisance maps are versioned and
validated;
\item theory discrepancy is separated by physical truncation layer and tested
against known higher levels;
\item prior predictive simulation rejects implausible or structurally invalid
regions before fitting;
\item synthetic recovery and \SBC\ validate the inference implementation;
\item exact and emulated calculations agree within a certified emulator
contract;
\item accepted samplers pass rank, tail, energy, mode, and Monte Carlo error
diagnostics;
\item one common posterior member identity propagates through every flavor,
parton, target, nuclear, scale, and process output;
\item covariance and nonlinear uncertainty interactions are retained;
\item family-wise predictive checks and grouped cross-validation are reported;
\item at least one nontrivial family is withheld from calibration before final
optimization;
\item the withheld family is evaluated under a frozen metric and decision rule;
\item failure triggers a typed physical revision rather than a TMD-specific
repair;
\item posterior predictive and uncertainty products are reproducible from
released manifests and hashes;
\item any experimental-design recommendation includes systematics, feasibility,
and factorization status;
\item no output is labeled ``predicted--withheld'' unless its normalization and
shape were excluded from calibration and the gate passed.
\end{enumerate}

\subsection{Program-level predictive gate}

The full DeuteronWigner model may be described as a genuinely predictive
microscopic model only when the acceptance criteria of Volumes~I--V and the
following inferential statement hold simultaneously:
\begin{quote}
A restricted, shared posterior over a declared microscopic Hamiltonian,
currents, Wilson dynamics, matched nuclear operators, LF-to-QCD matching, and
controlled discrepancies predicts at least one nontrivial nucleon or deuteron
observable family that was excluded from calibration, with correlated
uncertainties and without a family-specific normalization or shape parameter.
\end{quote}

A preferred first gate is:
\begin{equation}
 \caldata
 =
 \left\{
 \text{state/current data},
 \text{selected unpolarized/helicity data},
 \text{restricted Sivers information}
 \right\},
 \label{eq:first-calibration-gate}
\end{equation}
followed by one or more withheld predictions from:
\begin{equation}
 \holddata
 \in
 \left\{
 \begin{gathered}
 \text{Boer--Mulders},\quad
 \text{tensor naive-\(T\)-odd},\quad
 \text{antiquark phase},\\
 \text{gluon \(f/d\) hierarchy},\quad
 \text{\(b_1\) or tagged tensor family}
 \end{gathered}
 \right\}.
 \label{eq:first-holdout-gate}
\end{equation}
The exact first gate depends on data independence, scheme matching, and process
factorization available at implementation time.

\subsection{Handoff to implementation}

The next software prompt after C5 should not be written until the C5 report
identifies:
\begin{itemize}
\item actual Wilson-path, cut-ledger, phase-budget, and link-odd result APIs;
\item the retained validation-only state and common-overlap interfaces;
\item the final provenance and isolation objects;
\item exact matching-status and unresolved-physics fields;
\item computational cost and differentiability of the pilot.
\end{itemize}
The first inference coding package should then implement only I0--I1:
typed parameter, dataset-ancestry, and likelihood objects, together with
analytic validation.
It must not fit the accepted phenomenological boundary or promote the C5 pilot
to physical data analysis.

\appendix

\section{Stable formal requirements}
\label{app:requirements}

\begin{longtable}{L{0.16\textwidth}L{0.74\textwidth}}
\toprule
ID & Requirement\\
\midrule
\endhead
V6.OWN &
Every inferred quantity has one declared physical, nuisance, or discrepancy
owner; named-TMD-only free parameters are rejected.\\
V6.GRAPH &
The probabilistic graph composes only typed deterministic and conditional maps.\\
V6.DATA &
Every dataset has observable, convention, scheme, covariance, ancestry, and
partition identity.\\
V6.NODOUBLE &
Statistically overlapping data products cannot enter as independent
likelihood terms.\\
V6.SPLIT &
Calibration, diagnostic, and withheld partitions are frozen before final
optimization.\\
V6.LIKE &
Each likelihood is normalized, typed, and validated with its nuisance and
Jacobian structure.\\
V6.PRIOR &
Priors encode physical scales and uncertainty, while exact constraints remain
structural.\\
V6.PPC0 &
Prior predictive simulation is required before calibration.\\
V6.DISC &
Discrepancy is decomposed by physical source, power counting, support, and
correlation.\\
V6.IDENT &
Physical parameter and discrepancy identifiability is evaluated by Jacobians,
profiles, posterior correlations, and cross-channel information.\\
V6.EMUL &
Every emulator carries domain, covariance, derivative, exact-validation, and
fallback information.\\
V6.SYNTH &
Synthetic parameter recovery is required for each activated inference layer.\\
V6.SBC &
Simulation-based calibration validates the implemented posterior algorithm on
tractable instances.\\
V6.MCMC &
Accepted posterior runs pass rank, tail, energy, mode, and Monte Carlo error
diagnostics.\\
V6.MEMBER &
One posterior member remains indivisible across every physical output.\\
V6.COV &
Cross-observable covariance and nonlinear uncertainty interactions are
retained.\\
V6.PRED &
Latent-physics and future-measurement predictive distributions remain
distinct.\\
V6.PPC &
Posterior predictive checks include physical residual structure, not only one
global statistic.\\
V6.CV &
Cross-validation units follow conditional independence and data ancestry.\\
V6.ROBUST &
Prior, discrepancy, emulator, numerical, and formal-choice sensitivity is
reported.\\
V6.HOLD &
At least one nontrivial family is protected from calibration by a frozen
validation record.\\
V6.UNBLIND &
Unblinding is immutable and any subsequent revision creates a new model
version and holdout.\\
V6.REVISE &
A failed withheld gate triggers a typed physical revision; curve-specific
repair is forbidden.\\
V6.REPORT &
Members, covariance, evidence status, manifests, and hashes accompany every
release.\\
V6.DESIGN &
Experimental design optimizes information subject to cost, systematics, and
factorization status.\\
\bottomrule
\end{longtable}

\section{Minimal parameter-ownership schema}
\label{app:parameter-schema}

A machine-readable parameter record contains:
\begin{verbatim}
{
  "parameter_id": "wilson.eikonal.coupling.shared",
  "owner_type": "WILSON_KERNEL",
  "owner_object_id": "rescattering_kernel.v1",
  "role": "PHYSICAL",
  "units": "dimensionless",
  "support": {"transform": "log", "lower": 0.0},
  "symmetry_class": ["color_fundamental", "link_oriented"],
  "regulator_scope": ["uv_regulator_id", "ir_regulator_id"],
  "affected_operator_blocks": ["quark_sivers", "quark_boer_mulders"],
  "forbidden_direct_attachments": ["named_tmd_curve"],
  "prior_id": "prior.wilson_coupling.v1",
  "differentiability": "AD",
  "version": 1
}
\end{verbatim}

\section{Minimal dataset and partition schema}
\label{app:data-schema}

\begin{verbatim}
{
  "dataset_id": "experiment.observable.release.v1",
  "observable_id": "typed_process_observable",
  "experiment_id": "experiment",
  "event_ancestry_ids": ["raw_event_block_hash"],
  "kinematics": {"variables": [], "bin_edges": []},
  "units": "declared",
  "conventions": ["spin", "azimuth", "tensor", "normalization"],
  "scheme": {"uv": null, "rapidity": null, "factorization": null},
  "covariance_id": "covariance.full.v1",
  "likelihood_id": "likelihood.family.v1",
  "nuisance_ids": [],
  "partition": "CALIBRATION | DIAGNOSTIC | WITHHELD",
  "embargo_token": null,
  "source_hashes": [],
  "version": 1
}
\end{verbatim}

\section{Minimal posterior-member schema}
\label{app:member-schema}

\begin{verbatim}
{
  "member_id": "posterior.chain.draw",
  "model_label": "declared_discrete_variant",
  "parameter_values": {"parameter_id": "value"},
  "nuisance_values": {},
  "discrepancy_member_ids": {},
  "state_member_ids": {"nucleon": "", "deuteron": ""},
  "operator_member_ids": {"gtmd": "", "wilson": "", "nuclear": ""},
  "matching_evolution_ids": {"matching": "", "evolution": ""},
  "log_density": {"prior": 0.0, "likelihood_by_family": {}, "posterior": 0.0},
  "sampler_diagnostics": {},
  "artifact_hashes": {},
  "evidence_status": "CALIBRATED",
  "version": 1
}
\end{verbatim}

\section{Minimal withheld-validation schema}
\label{app:validation-schema}

\begin{verbatim}
{
  "validation_id": "holdout.observable_family.v1",
  "model_hash": "",
  "posterior_hash": "",
  "family": "withheld_family",
  "data_ancestry_ids": [],
  "prediction_artifact_hash": "",
  "metrics": [
    {"metric": "LOG_SCORE", "range": []},
    {"metric": "COVERAGE", "level": 0.68, "range": []}
  ],
  "structural_gates": [],
  "unblinding": {"time": null, "token_hash": ""},
  "decision": "FROZEN | PASS | PARTIAL | FAIL",
  "permitted_revision_classes": [],
  "replacement_holdout_required": true,
  "version": 1
}
\end{verbatim}

\section{Uncertainty-component dictionary}
\label{app:uncertainty-dictionary}

\begin{longtable}{L{0.28\textwidth}L{0.62\textwidth}}
\toprule
Component & Required identity\\
\midrule
\endhead
Calibration/statistical &
Posterior variation of shared physical parameters under the declared data
likelihood.\\
Experimental systematic &
Explicit nuisance parameter or covariance block with experiment and ancestry
identity.\\
Hamiltonian EFT &
Expansion order, expansion parameter, breakdown scale, coefficient prior, and
correlation kernel.\\
Fock/basis &
Retained sectors, basis level, matched comparison map, and convergence model.\\
Regulator/renormalization &
Regulator family, refitting protocol, renormalization conditions, and matched
observable residual.\\
Wilson/eikonal &
Wilson order, intermediate-state space, pole/cut convention, regulator, and
link/color identity.\\
Matching &
Closed operator basis, matching order, scale, scheme, and omitted-operator
remainder.\\
TMD evolution &
Anomalous-dimension order, Collins--Soper kernel, large-\(b\) model, path, and
threshold choices.\\
Nuclear many-body &
Nuclear sectors, interaction/current order, coherent amplitudes, and
partonic--nuclear subtraction.\\
Process/power &
Factorization status, hard/fragmentation order, \(W+Y\), power corrections,
and Glauber assessment.\\
Emulator &
Training design, domain, validation residual, covariance, and derivative
accuracy.\\
Numerical &
Precision, quadrature, grid, interpolation, solver, transform, and Monte Carlo
error.\\
\bottomrule
\end{longtable}

\section{Compact acceptance checklist}
\label{app:checklist}

Before any result is described as a prediction, answer:
\begin{enumerate}
\item Which microscopic parameters generate it?
\item Which calibration data constrain those parameters?
\item Is the predicted family absent from the calibration likelihood and its
data ancestry?
\item Which exact, truncation, matching, and process gates apply?
\item Does every posterior member produce the result through the same typed
forward chain?
\item Are cross-observable covariance and non-Gaussian tails retained?
\item Which uncertainty sources are probabilistic, which are envelopes, and
which are unresolved?
\item Was the prediction artifact frozen before unblinding?
\item Did the predeclared gate pass without post-unblinding adjustment?
\item Which future observation could falsify the retained claim?
\end{enumerate}


\small
\begin{thebibliography}{99}

\bibitem{GelmanWorkflow2020}
A.~Gelman, A.~Vehtari, D.~Simpson, C.~C.~Margossian, B.~Carpenter,
Y.~Yao, L.~Kennedy, J.~Gabry, P.-C.~B\"urkner, and M.~Modr\'ak,
``Bayesian workflow,''
arXiv:2011.01808.

\bibitem{Wesolowski2016}
S.~Wesolowski, N.~Klco, R.~J.~Furnstahl, D.~R.~Phillips, and A.~Thapaliya,
``Bayesian parameter estimation for effective field theories,''
J.\ Phys.\ G \textbf{43}, 074001 (2016), arXiv:1511.03618.

\bibitem{Furnstahl2015}
R.~J.~Furnstahl, N.~Klco, D.~R.~Phillips, and S.~Wesolowski,
``Quantifying truncation errors in effective field theory,''
Phys.\ Rev.\ C \textbf{92}, 024005 (2015), arXiv:1506.01343.

\bibitem{Melendez2019}
J.~A.~Melendez, R.~J.~Furnstahl, D.~R.~Phillips, M.~T.~Pratola, and
S.~Wesolowski,
``Quantifying correlated truncation errors in effective field theory,''
Phys.\ Rev.\ C \textbf{100}, 044001 (2019), arXiv:1904.10581.

\bibitem{KennedyOHagan2001}
M.~C.~Kennedy and A.~O'Hagan,
``Bayesian calibration of computer models,''
J.\ R.\ Stat.\ Soc.\ B \textbf{63}, 425 (2001),
doi:10.1111/1467-9868.00294.

\bibitem{BrynjarsdottirOHagan2014}
J.~Brynjarsd\'ottir and A.~O'Hagan,
``Learning about physical parameters: the importance of model discrepancy,''
Inverse Problems \textbf{30}, 114007 (2014),
doi:10.1088/0266-5611/30/11/114007.

\bibitem{Plumlee2017}
M.~Plumlee,
``Bayesian calibration of inexact computer models,''
J.\ Am.\ Stat.\ Assoc.\ \textbf{112}, 1274 (2017),
doi:10.1080/01621459.2016.1211016.

\bibitem{HoffmanGelman2014}
M.~D.~Hoffman and A.~Gelman,
``The No-U-Turn sampler: adaptively setting path lengths in Hamiltonian Monte
Carlo,''
J.\ Mach.\ Learn.\ Res.\ \textbf{15}, 1593 (2014), arXiv:1111.4246.

\bibitem{Talts2018}
S.~Talts, M.~Betancourt, D.~Simpson, A.~Vehtari, and A.~Gelman,
``Validating Bayesian inference algorithms with simulation-based
calibration,''
arXiv:1804.06788.

\bibitem{VehtariRhat2021}
A.~Vehtari, A.~Gelman, D.~Simpson, B.~Carpenter, and P.-C.~B\"urkner,
``Rank-normalization, folding, and localization: an improved
\(\widehat R\) for assessing convergence of MCMC,''
Bayesian Anal.\ \textbf{16}, 667 (2021), arXiv:1903.08008.

\bibitem{VehtariLOO2017}
A.~Vehtari, A.~Gelman, and J.~Gabry,
``Practical Bayesian model evaluation using leave-one-out cross-validation
and WAIC,''
Stat.\ Comput.\ \textbf{27}, 1413 (2017), arXiv:1507.04544.

\bibitem{VehtariPSIS2024}
A.~Vehtari, D.~Simpson, A.~Gelman, Y.~Yao, and J.~Gabry,
``Pareto smoothed importance sampling,''
J.\ Mach.\ Learn.\ Res.\ \textbf{25}, 72:1 (2024), arXiv:1507.02646.

\bibitem{YaoStacking2018}
Y.~Yao, A.~Vehtari, D.~Simpson, and A.~Gelman,
``Using stacking to average Bayesian predictive distributions,''
Bayesian Anal.\ \textbf{13}, 917 (2018), arXiv:1704.02030.

\bibitem{GneitingRaftery2007}
T.~Gneiting and A.~E.~Raftery,
``Strictly proper scoring rules, prediction, and estimation,''
J.\ Am.\ Stat.\ Assoc.\ \textbf{102}, 359 (2007),
doi:10.1198/016214506000001437.

\bibitem{GneitingCalibration2007}
T.~Gneiting, F.~Balabdaoui, and A.~E.~Raftery,
``Probabilistic forecasts, calibration and sharpness,''
J.\ R.\ Stat.\ Soc.\ B \textbf{69}, 243 (2007),
doi:10.1111/j.1467-9868.2007.00587.x.

\bibitem{Frame2018}
D.~Frame, R.~He, I.~Ipsen, D.~Lee, D.~Lee, and E.~Rrapaj,
``Eigenvector continuation with subspace learning,''
Phys.\ Rev.\ Lett.\ \textbf{121}, 032501 (2018), arXiv:1711.07090.

\bibitem{ZhangFurnstahl2022}
X.~Zhang and R.~J.~Furnstahl,
``Fast emulation of quantum three-body scattering,''
Phys.\ Rev.\ C \textbf{105}, 064004 (2022), arXiv:2110.04269.

\bibitem{Drischler2021}
C.~Drischler, M.~Qui\~nonez, P.~G.~Giuliani, A.~E.~Lovell, and F.~M.~Nunes,
``Toward emulating nuclear reactions using eigenvector continuation,''
Phys.\ Lett.\ B \textbf{823}, 136777 (2021), arXiv:2108.08269.

\bibitem{ChalonerVerdinelli1995}
K.~Chaloner and I.~Verdinelli,
``Bayesian experimental design: a review,''
Stat.\ Sci.\ \textbf{10}, 273 (1995).

\bibitem{CoonsHuan2025}
T.~E.~Coons and X.~Huan,
``A multi-fidelity estimator of the expected information gain for Bayesian
optimal experimental design,''
arXiv:2501.10845.

\end{thebibliography}

\end{document}
